% concatenated from simatos.tex
\documentclass[]{article}

\RequirePackage{amsthm,amsmath,amsfonts,amssymb}
\RequirePackage[numbers]{natbib}
\RequirePackage{graphicx}
\usepackage{algorithm}
\usepackage[noend]{algpseudocode}
\usepackage{color, dsfont}
\usepackage{subcaption}
\renewcommand{\d}{\textrm{d}}
\newcommand{\E}{\mathbb{E}}
\newcommand{\nfin}{{n_p}}
\newcommand{\npar}{{n}}
\newcommand{\Indicator}[1]{\mathds{1}\left(#1\right)}
\renewcommand{\P}{\mathbb{P}}
\newcommand{\R}{\mathbb{R}}
\newcommand{\Var}{\mathbb{V}\text{\textnormal{ar}}}
\newcommand{\Cov}{\mathbb{C}\text{\textnormal{ov}}}
\newcommand{\tr}{\text{\textnormal{tr}}}
\newcommand{\Tr}{\text{\rm Tr}}
\newcommand{\Proj}{{\textnormal{\textrm{Proj}}}} %{\textnormal{\textrm{proj}}}
\newcommand{\proj}{{\textnormal{\textrm{proj}}}} %{\textnormal{\textrm{proj}}}
\newcommand{\Lone}{\stackrel{\text{\textnormal{L\textsubscript{1}}}}{\longrightarrow}}
\newcommand{\lone}{\stackrel{\text{\textnormal{L\textsubscript{1}}}}{\to}}
\newcommand{\diag}{{\rm diag}}
\newcommand{\op}{{\rm op}}
\newcommand{\lin}{{\rm lin}}
\newcommand{\quadratic}{{\rm quad}}
\newcommand{\fin}{{\rm fin}}
\newcommand{\one}{{\bf 1}}
\newcommand{\iCE}{{\textnormal{\textrm{iCE}}}} %{\textnormal{\textrm{proj}}}
\newcommand{\correction}[1]{{\color{red}#1}}

%%%%%%%%%%%%%%%%%%%%%%%%%%%%%%%%%%%%%%%%%%%%%%
%%                                          %%
%% For Axiom, Claim, Corollary, Hypothesis, %%
%% Lemma, Theorem, Proposition              %%
%% use \theoremstyle{plain}                 %%
%%                                          %%
%%%%%%%%%%%%%%%%%%%%%%%%%%%%%%%%%%%%%%%%%%%%%%
\theoremstyle{plain}
\newtheorem{lemma}{Lemma}[section]
\newtheorem{prop}[lemma]{Proposition}
\newtheorem{corollary}[lemma]{Corollary}
\newtheorem{thm}[lemma]{Theorem}
%%%%%%%%%%%%%%%%%%%%%%%%%%%%%%%%%%%%%%%%%%%%%%
%%                                          %%
%% For Assumption, Definition, Example,     %%
%% Notation, Property, Remark, Fact         %%
%% use \theoremstyle{definition}            %%
%%                                          %%
%%%%%%%%%%%%%%%%%%%%%%%%%%%%%%%%%%%%%%%%%%%%%%
\theoremstyle{definition}
\newtheorem{rk}[lemma]{Remark}
\newtheorem{assumption}{Assumption}

\begin{document}

\title{Phase transition for conditional covariance matrices estimated by importance sampling, and implications for cross-entropy schemes in high dimension}

\author{J.\ Beh$^{1,2}$ \and J.\ Morio$^{1}$ \and F.\ Simatos$^{2}$}
\date{%
    $^1${\small \it ONERA/DTIS, Université de Toulouse, F-31055 Toulouse}\\%
    $^2${\small \it Fédération ENAC ISAE-SUPAERO ONERA, Université de Toulouse, 31000 Toulouse}\\[2ex]%
    \today
}


\maketitle

\begin{abstract}
Motivated by the estimation of covariance matrices by importance sampling arising in the cross-entropy (CE) algorithm, we study a random matrix model $\hat \Sigma = {\bf X} L {\bf X}^\top$ with two distinct features: $\bf X$ and $L$ are dependent, and $L$ is heavy-tailed. In the high-dimensional regime $d \to \infty$, we prove under suitable assumptions that a phase transition occurs in the polynomial regime $n = d^\kappa$, with $n$ the sample size. Namely, we prove that $\lVert \hat \Sigma - \E \hat \Sigma \rVert_\op \Rightarrow 0$ if and only if $\kappa > \kappa_*$ for some threshold $\kappa_*$ determined by the behavior of the maximum likelihood ratios. Moreover, we identify general situations where $\kappa_* = 1/\lambda_1$, with $\lambda_1$ the smallest eigenvalue of the covariance matrix of the auxiliary distribution used to estimate $\hat \Sigma$ by importance sampling. This suggests that importance sampling will work better with covariance matrices having a large smallest eigenvalue. We carry this insight into recent CE schemes proposed to estimate the probability of high-dimensional rare events. Through numerical simulations, we demonstrate that better CE schemes are also the ones with larger smallest eigenvalue, even though these algorithms were not designed to smooth the spectrum. This new spectral interpretation raises stimulating questions and opens research directions for the design of efficient high-dimensional algorithms.
\end{abstract}


\section{Introduction}
\subsection{Motivation: cross-entropy in high dimension}

Consider the problem of estimating the probability $p = \P(X \in A)$ of an event $A \subset \R^d$. In the reliability context where $p$ is small, i.e., $A$ is rare, the naive Monte Carlo method is not efficient and one has to resort to better numerical schemes. for instance adaptive splitting methods such as subset simulation or adaptive importance sampling methods such as cross-entropy. Importance sampling is a general technique widely used in this context. It estimates $p$ from $\nfin$ i.i.d.\ samples $X_i$ drawn according to an auxiliary distribution $g$ in the following way:
\begin{equation} \label{eq:IS}
	\hat p = \frac{1}{\nfin} \sum_{k=1}^\nfin \frac{f(X_i)}{g(X_i)} \xi_A(X_i)
\end{equation}
where $\xi_A$ is the indicator function of the set $A$ and $f$ is the original density of interest, so that $p = \int_{\R^d} \xi_A f$. 

The choice of the auxiliary distribution~$g$ is of uttermost importance for the efficiency of the importance sampling scheme. Here we are specifically interested in the behavior of the prominent cross-entropy (CE) algorithm, which is an iterative algorithm that outputs a family of auxiliary distributions $(\hat g_t, t = 0, 1, \dots)$. CE works for rare events of the form $A = \{x: \varphi(x) \geq 0\}$ for some measurable function $\varphi: \R^d \to \R$, so that the sought probability can be written as $p = \P(\varphi(X) \geq 0)$. The idea is to sequentially approach the set $\{\varphi \geq 0\}$ by simpler sets $\{\varphi \geq q_t\}$, with probabilities $\P(\varphi(X) \geq q_t)$ approaching $p$. A deterministic version of CE is presented in Algorithm~\ref{alg:CE-det} showing how the $q_t$'s are iteratively computed: at iteration $t$ with current auxiliary distribution $g_t$, $q_t$ is computed as the $\rho$-quantile of $\varphi(X)$ under~$g_t$, and the next auxiliary distribution $g_{t+1}$ is then the Gaussian distribution with mean and variance that of $f$ conditioned on $\varphi \geq q_t$. This deterministic version of CE involves the unknown parameters $g_t$, $q_t$, $A_t$, $\mu_t$ and $\Sigma_t$. The true CE scheme, presented in Algorithm~\ref{alg:CE}, is simply a stochastic version of Algorithm~\ref{alg:CE-det} where these unknown parameters are replaced by estimations $\hat g_t$, $\hat q_t$, $\hat A_t$, $\hat \mu_t$ and $\hat \Sigma_t$.
%
\begin{algorithm}[t]
\caption{One iteration of the deterministic version of CE} \label{alg:CE-det}
\begin{algorithmic}
\Require $\rho \in (0,1)$, current auxiliary distribution $g_t = N(\mu_t, \Sigma_t)$
\State 1. compute $q_t$ such that $\P_{g_t}(\varphi(X) \geq q_t) = \rho$;
\State 2. compute $\mu_{t+1}$ and $\Sigma_{t+1}$ the mean and variance of $f$ conditioned on $A_t = \{x: \varphi(x) \geq q_t\}$;
\State 3. iterate with $g_{t+1} = N(\mu_{t+1}, \Sigma_{t+1})$.
\end{algorithmic}
\end{algorithm}

\begin{algorithm}[t]
\caption{One iteration of CE} \label{alg:CE}
\begin{algorithmic}
\Require $\rho \in (0,1)$, current auxiliary distribution $\hat g_t = N(\hat \mu_t, \hat \Sigma_t)$, sample sizes~$m$ and $\npar$
\State 1. generate $Y_1, \ldots, Y_m$ i.i.d.\ according to $\hat g_t$, independently from the $Y$'s;
\State 2. rank the $Y$'s according to their values by $\varphi$: $\varphi(Y_{(1)}) \leq \cdots \leq \varphi(Y_{(m)})$ and define $\hat q_t = \varphi(Y_{(\lfloor (1-\rho) m\rfloor)})$ and $\hat A_t = \{x: \varphi(x) \geq \hat q_t\}$;
\State 3. generate $X_1, \ldots, X_\npar$ i.i.d.\ according to $\hat g_t$;
\State 4. compute
\begin{equation} \label{eq:Sigma-CE}
	\left\{ \begin{array}{rl}
		\hat \mu_{t+1} & \displaystyle = \frac{1}{\npar} \sum_{k=1}^{\npar} \ell(X_i) X_i\\
		\hat \Sigma_{t+1} & \displaystyle = \frac{1}{\npar} \sum_{k=1}^{\npar} \hat \ell(X_i) X_i X_i^\top - \hat \mu_{t+1} \hat \mu_{t+1}^\top
	\end{array} \right.
\end{equation}
where $\displaystyle \hat \ell = \frac{f \xi_{\hat A_t}}{\hat p_t \hat g_t}$ with $\displaystyle \hat p_t = \frac{1}{\npar} \sum_{k=1}^\npar \frac{f(X_i)}{\hat g_t(X_i)} \xi_{\hat A_t}(X_i)$;
\State 5. iterate with $\hat g_{t+1} = N(\hat \mu_{t+1}, \hat \Sigma_{t+1})$.
\end{algorithmic}
\end{algorithm}
%




We are interested in the behavior of this algorithm in the high-dimensional regime where $d \to \infty$, so all quantities of interest now implicitly rely on $d$. For instance, the original data $A$ and $f$ depend on the dimension, but this is not reflected in the notation to ease the readability. In high dimension, standard importance sampling techniques typically behave poorly. It is often observed in practice that the importance weights $\frac{f(X_i)}{g(X_i)}$ in~\eqref{eq:IS} degenerate, in the sense that one weight takes all the mass. In CE, the issue stems from the estimation~\eqref{eq:Sigma-CE} of the covariance matrix $\Sigma_t$ in step~$4$ of Algorithm~\ref{alg:CE}. Indeed, it is known that estimating high-dimensional covariance matrices is a very difficult problem~\cite{lam2019}, as it amounts to estimating a quadratic (in the dimension $d$) number of parameters. Thus, except if the sample size $\npar$ is large, the estimation $\hat \Sigma_t$ of $\Sigma_t$ will be noisy.

To overcome this problem, a recent stream of literature has proposed to reduce the dimension of the problem by projecting on subspaces of small dimensions. The general idea is to replace $\hat \Sigma_t$ in step 4 of Algorithm~\ref{alg:CE} by a covariance matrix that only updates variance terms in a small number of directions: typically, this amounts to replacing $\hat \Sigma_t$ by $\Proj_r(\hat \Sigma_t, {\bf v})$ where we define the following projection operator with ${\bf v} = (v_1, \dots, v_r)$ an orthonormal family:
\begin{equation} \label{eq:Proj}
	\Proj_r(\Sigma, {\bf v}) = \sum_{k=1}^r (\lambda_k - 1) v_k v_k^\top + I \ \text{ with } \ \lambda_k = v_k^\top \Sigma v_k.
\end{equation}
In words, $\Proj_r(\Sigma, {\bf v})$ keeps the same variance as $\Sigma$ in the directions $v_1, \ldots, v_r$ (the $\lambda_k$'s), while setting variance in the orthogonal subspace to the identity. The subscript $r$ refers to the dimension of the subspace which in practice is typically less than $3$. One should thus think of $\Proj_r(\Sigma, {\bf v})$ as a finite-rank perturbation of the identity, the so-called the spiked model. Algorithm~\ref{alg:CE-proj} presents a typical version of CE with projection: the main difference with Algorithm~\ref{alg:CE} is step~4, where the full covariance matrix $\hat \Sigma_{t+1}$ estimated by importance sampling is replaced by a simpler covariance matrix $\hat \Sigma^\proj_{t+1}$ obtained by projection. In particular, we see that in versions of CE with projection, the covariance matrices used in the intermediate Gaussian auxiliary densities $\hat g_t^\proj$ are finite-rank perturbations of the identity.
%
\begin{algorithm}[t]
\caption{One iteration of CE with projection}\label{alg:CE-proj}
\begin{algorithmic}
\Require $\rho \in (0,1)$, current auxiliary distribution $\hat g^\proj_t = N(\hat \mu^\proj_t, \hat \Sigma^\proj_t)$, sample sizes $m$ and $\npar$
\State 1. generate $Y_1, \ldots, Y_m$ i.i.d.\ according to $\hat g_t^\proj$ and define $\hat q_t$ and $\hat A_t$ from the $Y_i$'s similarly as in step 2 of Algorithm~\ref{alg:CE};
\State 2. generate $X_1, \ldots, X_\npar$ i.i.d.\ according to $\hat g_t^\proj$, independently from the $Y$'s, and compute $\hat p_t$, $\hat \mu_{t+1}$ and $\hat \Sigma_{t+1}$ from the $X_i$'s similarly as in step 4 of Algorithm~\ref{alg:CE} but with $\hat g^\proj_t$ instead of $\hat g_t$;
\State 3. compute the orthonormal family ${\bf v} = (v_1, \dots, v_r)$;
\State 4. define 
\begin{align}\label{eq:Sigma-proj}
\hat \mu^\proj_{t+1} = \hat \mu_{t+1} \text{ and } \hat \Sigma^\proj_{t+1} = \Proj_r(\hat \Sigma_{t+1}, {\bf v})
\end{align}
and iterate with $\hat g_{t+1}^\proj = N(\hat \mu^\proj_{t+1}, \hat \Sigma_{t+1}^\proj)$.
\end{algorithmic}
\end{algorithm}
%

Another motivation for considering spiked covariance matrices is that the rare event $A$ has often, in practice, a small intrinsic dimension, meaning that it can be described by a linear subspace of a much lower dimension than the dimension $d$ of the input. When $A = \{ \varphi \geq 0\}$ this means that $\varphi(x) = \varphi(Mx)$ for some $d \times d$ matrix $M$ with small rank (the intrinsic dimension). In this case, it is natural to only try and evaluate covariance terms in a small number of directions deemed influential. Such dimensionality-reduction behavior is widespread in signal processing, learning theory or data assimilation, to name only a few examples, see for instance the literature overview in~\cite{agapiou_importance_2017}.


Various propositions have been made for the projection directions ${\bf v}$, see for instance~\cite{el_masri2024, ELMASRI2021, uribe2021} and the discussion in~\cite[Section $2.3$]{beh2023insight}. In practice this choice has an important impact on the accuracy of the resulting estimator, but, to our knowledge, no theoretical result to date has assessed its influence on the performance of the final estimator $\hat p$.

\subsection{Estimation of covariance matrices by importance sampling}

As explained above, the high-dimensional performance of CE is strongly driven by the quality of the estimation of $\Sigma_t$ by importance sampling in~\eqref{eq:Sigma-CE}. Neglecting the rank-one term $\hat \mu_t \hat \mu_t^\top$, $\hat \Sigma_t$ follows a random matrix model that can be written as
\begin{equation} \label{eq:RMT}
	\hat \Sigma = \frac{1}{n} {\bf X} \hat L {\bf X}^\top
\end{equation}
with ${\bf X} = [X_1 \cdots X_\npar]$ the $\npar \times d$ matrix with columns $X_i$ and $\hat L$ the $n \times n$ diagonal matrix with entries $\hat \ell(X_i)$. This model, called the covariance matrix model, is quite standard in random matrix theory. It has mostly been studied assuming that $\hat L$ and $\bf X$ are independent with $\hat L$ light-tailed: under these assumptions, the relevant regime is where $n$ and $d$ grow proportionate, i.e., $d/n \to c$ for some $c \in (0,\infty)$~\cite{BAIK20061382, benaych-georges_eigenvalues_2011, capitaine2016spectrumdeformedrandommatrices, Paul_07}. However, these two assumptions do not hold here: $\hat L$ and $\bf X$ are not independent, and in high dimension, $\hat L$ is typically heavy-tailed. To the best of our knowledge, no existing result covers this case, and our main theoretical result makes progress in this direction. Namely, we show (see Theorem~\ref{thm:main} for a precise statement) that, when $A$ has a finite intrinsic dimension, an original phase transition occurs in the polynomial regime $n = d^\kappa$: for $\hat \Sigma$ as in~\eqref{eq:RMT}, the approximation $\hat \Sigma \approx \E(\hat \Sigma)$ is accurate if and only $\kappa > \kappa_*$ for some threshold $\kappa_*$ linked to the behavior of the maximum likelihood ratios $\hat \ell(X_i)$. From a technical standpoint, what allows us to control $\hat \Sigma - \E \hat \Sigma$ is that, under the spiked model and assuming that $A$ has a finite intrinsic dimension, $\hat L$ and $\bf X$ are only slightly dependent, and so we can essentially reduce the problem to the independent case where we can then use recent concentration results for the sum of independent random matrices~\cite{Brailovskaya24-0}.



\subsection{Paper organization}


The paper is organized as follows. Section~\ref{sec:main-thm} introduces the simplified random matrix model that serves as a proxy for covariance matrix estimation in CE schemes, and presents and discusses our main theoretical results. Section~\ref{sec:main-newlook} connects these results back to CE schemes through numerical experiments on classical benchmark examples. This provides a new perspective on the effect of projection strategies on their performance. Section~\ref{sec:conclu} outlines potential directions for future research opened by this new perspective. Finally, proofs are presented in Section~\ref{sec:proofs}.







\section{Phase transition for conditional covariance matrices estimated by importance sampling} \label{sec:main-thm}


\subsection{General notation}

In the rest of the paper, $N(\mu, \Sigma)$ denotes the Gaussian density with mean $\mu$ and covariance matrix $\Sigma$, with $I$ the identity matrix. For $f$ a density and $A \subset \R^d$ a measurable subset, $f|_A$ denotes the density $f$ conditioned on $A$: $f|_A = f \xi_A / p$ with $\xi_A$ the indicator of the set $A$ and $p = \P_f(X \in A)$. Here and in the sequel, when using probability or expectation operators, the subscript indicates the law of the corresponding variables, so for instance $\P_f(X \in A)$ means that $X$ is distributed according to $f$ (written $X \sim f$). When the law is not indicated as a subscript, the law of the relevant random variables is indicated in the text. For sequences of random variables, we denote by $\Rightarrow$ and $\Lone$ convergence in distribution and in L\textsubscript{1}, respectively.

For $W \subset \R^d$ a linear subspace, $P_W$ denotes the orthogonal projection on $W$ and $W_\perp$ denotes its orthogonal subspace. For a vector $x \in \R^d$, $\lVert x \rVert^2 = x^\top x = \langle x, x \rangle$ is the square L\textsubscript2 norm, with $\langle \cdot, \cdot \rangle$ the inner product. For a $d \times d$ matrix $M = (M_{ij})$, $\lVert \cdot \rVert_\op = \sup_{x: \lVert x \rVert = 1} \lvert x^\top M x \rvert$ denotes its operator norm, $\lvert M \rvert$ its determinant and $\Tr(M) = \sum_i M_{ii}$ its trace. When $M$ is symmetric with real eigenvalues, we denote by $\lambda_{\min}(M)$ and $\lambda_{\max}(M)$ its minimal and maximal eigenvalues, respectively. Note that all quantities depend on $d$ but this is omitted from the notation for clarity. 

\subsection{Main results}


In the rest of the paper, we consider $f = N(0,I)$ and $A \subset \R^d$ measurable with $p = \P_f(X \in A)$. Let $\Sigma_A$ be the covariance matrix of~$f$ conditioned on $A$:
\[ \Sigma_A = \E_f( Y Y^\top \mid Y \in A) - \mu_A \mu_A^\top \ \text{ with } \ \mu_A = \E_f(Y \mid Y \in A). \]
We are interested in the estimation of $\Sigma_A$ by importance sampling using as auxiliary distribution a centered Gaussian distribution whose covariance matrix is a finite-rank perturbation of the identity. More precisely, we will consider the following estimator of $\Sigma_A$:
\begin{equation} \label{eq:hat-Sigma}
	\hat \Sigma_A = \frac{1}{np} \sum_{i=1}^n \ell(X_i) \xi_A(X_i) X_i X_i^\top - \mu_A \mu_A^\top
\end{equation}
with $\ell = f/g$ and the $X_i$'s i.i.d.\ sampled according to $g = N(0, \Sigma)$. We assume that $\Sigma$ satisfies the following assumption, which states that it follows the spiked model, i.e., it is a finite-rank perturbation of the identity. Note that this assumption allows for $\Sigma = I$ by taking $r = \lambda_1 = 1$: this case corresponds to the standard Monte Carlo scheme where $g = f$.

\begin{assumption} \label{ass:Sigma}
	For each $d$, the covariance matrix $\Sigma$ is given by
	\begin{equation} \label{eq:Sigma}
		\Sigma = \sum_{k=1}^r (\lambda_k - 1) v_k v_k^\top + I
	\end{equation}
	for some $r \geq 1$, $0 < \lambda_1 \leq \lambda_2 \leq \cdots \leq \lambda_r$ and some orthonormal family $(v_1, \ldots, v_r)$, whose span is defined as $V = \text{span}(v_1, \ldots, v_r)$. Moreover, as $d$ varies we have that:
	\begin{itemize}
		\item $\sup_d r < \infty$;
		\item $\lambda_1 = \lambda_{\min}(\Sigma)$ is independent of $d$ and satisfies $0 < \lambda_1 \leq 1$;
		\item $\sup_d \lambda_{\max}(\Sigma) < \infty$.
	\end{itemize}
\end{assumption}


The estimator~\eqref{eq:hat-Sigma} is inspired by the estimation step~\eqref{eq:Sigma-proj} in CE with projection. The main differences are that 1/ we assume in~\eqref{eq:hat-Sigma} that the probability $p$ and the mean $\mu_A$ are known, whereas they are estimated in~\eqref{eq:Sigma-proj}, and 2/ the auxiliary distribution is not centered in~\eqref{eq:Sigma-proj}. The simplified model~\eqref{eq:hat-Sigma} allows to focus on the main difficulty in~\eqref{eq:Sigma-proj}, and our results could probably be extended with estimations of $p$ and $\mu_A$ and non-centered $g$ at the expense of more involved technical derivations but with no additional insight.

Our second assumption concerns the rare event of interest $A$: for reasons outlined in the introduction, we assume that it has a finite intrinsic dimension.

\begin{assumption} \label{ass:FID}
	For each $d$, there exists a linear subspace $U \subset \R^d$ such that
	\begin{equation} \label{eq:FID}
		\forall x \in \R^d, \ x \in A \Leftrightarrow P_{U} x \in A.
	\end{equation}
	Moreover, as $d$ varies we have $\sup_d \dim(U) < \infty$.
\end{assumption}

Thus, two subspaces play a key role:
\begin{itemize}
	\item $U$, which is enough to describe the set of interest $A$;
	\item and $V$, the set of directions in which variance terms are updated in the auxiliary distribution~$g$. 
\end{itemize}
Intuitively, we expect the estimation $\hat \Sigma_A \approx \Sigma_A$ to be better when $V$ is aligned with $U$. Indeed, Assumption~\ref{ass:FID} implies that $\Sigma_A x = x$ for $x \in U_\perp$ (see the forthcoming proof of Lemma~\ref{lemma:Sigma-A-mu-A-W}), which means that directions in the orthogonal of $U$ do not influence $\Sigma_A$. Thus, only the directions in $U$ are influential and so one should try to take the $v_k$'s belonging to $U$. 

However, although different choices for $\bf v$ have been proposed in the literature, the influence of this choice on the quality of the estimation has never been assessed theoretically. Here we make progress in this direction by studying the two extreme cases where either $V \subset U$ (``good'' case) or $V \subset U_\perp$ (``bad'' case). Our main result shows that, in both cases, a phase transition occurs in the polynomial regime $n = d^\kappa$: the estimation will be accurate if $\kappa > \kappa_*$ and not accurate if $\kappa < \kappa_*$. We show that the value of the threshold $\kappa_*$ is linked to the smallest eigenvalue $\lambda_1$ of~$\Sigma$ through the behavior of the maximum likelihood ratio appearing in~\eqref{eq:hat-Sigma}. To control the latter, we will make the following technical assumption.





\begin{assumption} \label{ass:MW}
	There exists $\gamma_* \geq 0$ such that:
	\begin{equation} \label{eq:MW}
		\frac{1}{n^\gamma} \max_{1 \leq i \leq n} \xi_A(X_i) \ell(X_i) \Rightarrow \left\{ \begin{array}{ll} 0 & \text{ if } \gamma > \gamma_*,\\
		+\infty & \text{ if } \gamma < \gamma_* \end{array} \right.
	\end{equation}
	where the $X_i$'s are i.i.d.\ drawn according to $g$. 
\end{assumption}

In the non-triangular case (no dependency on $d$), the fact that the maximum of $n$ i.i.d.\ random variables $Z_i$ grows polynomially is quite standard. More precisely, in the non-triangular case,~\eqref{eq:MW} holds with $\gamma_* = 0$ for light-tailed random variables, and when $Z_i$ is in the domain of attraction of an $\alpha$-stable distribution ($\alpha < 2$), then it holds with $\gamma_* = 1/\alpha$ (in this case, $\max_{i=1, \dots, n} Z_i$ grows like $n^{1/\alpha} s(n)$ for some slowly varying function~$s$). Thus, we assume that this standard behavior continues to hold in the triangular regime $d \to \infty$, which can be seen as a regularity assumption as the dimension increases. We can now state our main result.

\begin{thm} \label{thm:main}
	Assume that Assumptions~\ref{ass:Sigma},~\ref{ass:FID} and~\ref{ass:MW} hold, and let $\kappa_* = 1/(1-\gamma_*)$. Assume in addition that $\inf_d p > 0$ and that either $V \subset U$ or $V \subset U_\perp$. Then in the regime $n = d^\kappa$, the following phase transition holds:
	\begin{itemize}
		\item if $\kappa > \kappa_*$, then $\lVert \hat \Sigma_A - \Sigma_A \rVert_\op \Rightarrow 0$;
		\item if $\kappa < \kappa_*$, then $\lVert \hat \Sigma_A - \Sigma_A \rVert_\op \Rightarrow \infty$ and more precisely, we have $\lambda_{\max}(\hat \Sigma_A) \Rightarrow \infty$ while $\sup_d \lambda_{\max}(\Sigma_A) < \infty$.
	\end{itemize}
	Moreover, if $V \subset U_\perp$ then $\kappa_* = 1/\lambda_1$ while if $V \subset U$, then $1 \leq \kappa_* \leq 1/\lambda_1$.
\end{thm}

The assumption $\inf_d p > 0$ is crucial for our results. This assumption is common in previous work studying the high-dimensional behavior of importance sampling, see the discussion in~\cite[Section $2.4$]{beh2023insight}. To our knowledge, the case $p \to 0$ is essentially unknown, and widely different behavior requiring significant different arguments could arise. Importantly, it was proved in~\cite[Corollary $4.2$]{beh2023insight} that the assumption $\inf_d p > 0$ implies that
\begin{equation} \label{eq:inf-p}
	\sup_d \lVert \mu_A \rVert < \infty \ \text{ and } \ 0 < \inf_d \lambda_{\min}(\Sigma_A) \leq \sup_d \lambda_{\max}(\Sigma_A) < \infty.
\end{equation}


The next result shows that the range $[1,1/\lambda_1]$ for $\kappa_*$ in the case $V \subset U$ is sharp, i.e., every value can be attained.


\begin{prop} \label{prop:range}
	Assume that $A$ is of the form $A = \{x \in \R^d: \lvert \langle u, x \rangle \rvert \leq K \}$ for some $u \in \R^d$ with $\lVert u \rVert = 1$ and $K > 0$. Then Assumption~\ref{ass:FID} is satisfied with $U = {\rm span}(u)$, and for any fixed $\lambda_1 \in (0,1)$, Assumption~\ref{ass:Sigma} is also satisfied if $V = U$. Moreover, if $K = 1 + \sqrt{2 \alpha \lambda_1 \log n}$ for some $\alpha \in (0,1]$, then $\inf_d p > 0$ and Assumption~\ref{ass:MW} is satisfied with $\gamma_* = \alpha (1-\lambda_1)$.
\end{prop}

Of course, this result is very artificial as we have $p \to 1$, nonetheless it serves its purpose of showing that every value between $1$ and $1/\lambda_1$ is possible for $\kappa_*$. A slightly more interesting example is when $A = \{x: \langle x, u \rangle \geq K\}$ for some fixed $K \in \R$, not depending on $d$. In this case, even in the ``good'' case $V = {\rm span}(u)$, it is not hard to prove that $\kappa_* = 1/\lambda_1$. In particular, it is striking to realize that, depending on the set $A$, the good choice $V \subset U$ may not improve the performance of the bad choice $V \subset U_\perp$, with both needing $n \gg d^{1/\lambda_1}$ to behave correctly.




\section{Efficiency of recent CE schemes: a new spectral interpretation} \label{sec:main-newlook}

\subsection{Implications of Theorem~\ref{thm:main} for CE schemes with projection} \label{sub:implications}

Theorem~\ref{thm:main} has direct implications for CE schemes with projection described in Algorithm~\ref{alg:CE-proj}. Indeed, step~2 of Algorithm~\ref{alg:CE-proj} involves the estimation of the covariance matrix $\Sigma_{t+1} = \Sigma_{A_t}$ by importance sampling with auxiliary distribution $\hat g^\proj_t = N(\hat \mu^\proj_t, \hat \Sigma^\proj_t)$: since $\hat \Sigma^\proj_t$ is obtained by projection~\eqref{eq:Proj}, by construction it follows the spiked model of Assumption~\ref{ass:Sigma}. In other words, Theorem~\ref{thm:main} sheds light on CE schemes with projection when considering $A = A_t$ and $\Sigma = \hat \Sigma^\proj_t$. In particular, the prime importance of $\lambda_1 = \lambda_{\min}(\Sigma)$ discussed after Theorem~\ref{thm:main} suggests that Algorithm~\ref{alg:CE-proj} will behave better if $\lambda_{\min}(\hat \Sigma^\proj_t)$ is larger: indeed, a small smallest eigenvalue translates to a large sample size needed to accurately learn the target covariance matrix. To a lesser extent, the assumption $\sup \lambda_{\max}(\Sigma) < \infty$ in Assumption~\ref{ass:Sigma} also suggests that a small $\lambda_{\max}(\hat \Sigma^\proj_t)$ is beneficial. But in practice, projection schemes almost always only update variance terms in directions of small variance, leading to $\lambda_{\max}(\hat \Sigma^\proj_t) = 1$.

On this background, several projection directions $\bf v$ have recently been proposed. However, the rationale for this projection step is not to smooth the spectrum of the auxiliary covariance matrix, but rather to decrease the number of parameters to be estimated, thereby reducing noise while focusing on relevant parameters. Projection directions are then chosen according to some optimality criterion, typically minimizing the Kullback--Leibler divergence with $f|_A$.

In the following, we take a new look at these algorithms and pay special attention to the influence of the choice of projection directions on the spectrum of the estimated covariance matrices. Numerical results presented below show that projection directions that work well are also the ones with the largest smallest eigenvalue $\lambda_{\min}(\hat \Sigma^\proj_t)$. This is quite surprising since, as explained above, these algorithms were not designed to act on the spectrum. Nonetheless, this gives credit to our main result, Theorem~\ref{thm:main}, for the understanding of the performance of importance sampling schemes, and opens the way for the design of efficient importance sampling schemes in high dimensions.






\subsection{Spectral behavior of CE schemes}

\subsubsection{Numerical set-up}

%
\begin{table}[t]
	\begin{center}
	\begin{tabular}{|c|c|c|c|}
		\hline
		Test function & $\varphi_\lin$ & $\varphi_\quadratic$ & $\varphi_\fin$ \\\hline
		Dimension $d$ & $100$ & $334$ & $334$ \\\hline
		Probability $p$ & $2.9 \cdot 10^{-7}$ & $6.6 \cdot 10^{-6}$ & $1.8 \cdot 10^{-6}$ \\\hline
		Learning sample size $n$ & $10000$ & $5000$ & $5000$ \\\hline
		Final sample size $\nfin$ & \multicolumn{3}{c|}{$2000$} \\\hline
		Repetition $N$ & \multicolumn{3}{c|}{$200$} \\\hline
	\end{tabular}
	\end{center}
	\caption{Parameters used for the numerical simulations depending on the test function considered.}
	\label{tab:parameters}
\end{table}
%


To check numerically that efficient CE schemes have large $\lambda_{\min}(\hat \Sigma^\proj_t)$, we consider two CE schemes, two choices of projection directions and three functions $\varphi$ (recall that, in CE schemes, rare events are of the form $A = \{x \in \R^d: \varphi(x) \geq 0\}$).

% 
\begin{algorithm}[t]
\caption{One iteration of improved CE (iCE) with projection}\label{alg:iCE-proj}
\begin{algorithmic}
\Require $\delta > 0$, current auxiliary distribution $\hat g^\iCE_t = N(\hat \mu^\proj_t, \hat \Sigma^\proj_t)$, $\hat \sigma_t$, sample sizes $m$ and $\npar$
\State 1. generate $Y_1, \ldots, Y_m$ i.i.d.\ according to $\hat g^\iCE_t$ and compute
\[ \hat \sigma_{t+1} = \arg \min_{\sigma \in (0, \hat \sigma_t)} (\hat \delta_t(\sigma) - \delta )^2, \]
where $\hat \delta_t(\sigma)$ is the empirical coefficient of variation of $F_N(\varphi(Y_i)/\sigma) \frac{f(Y_i)}{\hat g^\iCE_t(Y_i)}$:
\begin{equation}\label{eq:iCE-delta}
    \hat \delta_t(\sigma) = \sqrt{ m \sum_{i=1}^m \left( \frac{f(Y_i)}{\hat g^\iCE_t(Y_i)} F_N\left( \frac{\varphi(Y_i)}{\sigma} \right) \right)^2} \biggl/ \sum_{i=1}^m \frac{f(Y_i)}{\hat g^\iCE_t(Y_i)} F_N \left( \frac{\varphi(Y_i)}{\sigma} \right);
\end{equation}
\State 2. generate $X_1, \ldots, X_\npar$ i.i.d.\ according to $\hat g^\iCE_t$, independently from the $Y_k$'s, and compute $\hat \mu_{t+1}$ and $\hat \Sigma_{t+1}$ from the $X$'s similarly as in~\eqref{eq:Sigma} with
\[ \ell = \frac{f F_N(\varphi(\cdot)/\hat \sigma_{t+1})}{\hat E \hat g^\iCE_t} \ \text{ where } \hat E = \frac{1}{\npar} \sum_{i=1}^\npar \frac{f(X_i)}{\hat g^\iCE_t(X_i)} F_N(\varphi(X_i) / \hat \sigma_{t+1}); \]
\State 3. define $\hat \mu^\proj_{t+1} = \hat \mu_{t+1}$ and compute $ \hat \Sigma^\proj_{t+1}$ from $\hat \Sigma_{t+1}$ similarly as in step 3 of Algorithm~\ref{alg:CE-proj};
\State 4. iterate with $\hat g^\iCE_{t+1} = N(\hat \mu^\iCE_{t+1}, \hat \Sigma^\proj_{t+1})$.
\end{algorithmic}
\end{algorithm}
% 
\begin{description}
	\item[CE schemes:] we consider two CE schemes: the standard CE scheme (Algorithm~\ref{alg:CE}) and also the improved CE scheme (iCE) proposed in~\cite{papaioannou_improved_2019}. The complete description of iCE (with projection) is provided in Algorithm~\ref{alg:iCE-proj} but the main idea is quite simple: it amounts to replacing the indicator~$\xi_A$ by a smooth approximation. This is the purpose of the terms $F_N(\varphi(\cdot) / \sigma)$ in Algorithm~\ref{alg:iCE-proj}, which satisfies $F_N(\varphi(\cdot) / \sigma) \approx \xi_A$ for small $\sigma$'s;
	\item[Projection directions:] we will compare CE and iCE without projection and with projection for two choices of projection directions. Each time, we consider for simplicity only one direction ($r = 1$). During iteration $t$, in step~3 of Algorithms~\ref{alg:CE-proj} and~\ref{alg:iCE-proj}, we will consider projecting on 1/ the eigenvector of $\hat \Sigma_{t+1}$ associated to its smallest eigenvalue and 2/ the current mean $\hat \mu_{t+1}$. The first choice is a simplified version of the algorithm proposed in~~\cite{el_masri2024}, while the second one was proposed in~\cite{ELMASRI2021}. This leads to three variants of CE and iCE denoted by CE, CE-eig and CE-mean for CE and iCE, iCE-eig and iCE-mean for iCE;
	\item[Functions $\varphi$:] we will test these six CE schemes on three different test functions, namely a linear function $\varphi_{\lin}(x) = x^\top \one - 5$ with $\one = \frac{1}{\sqrt d} (1, \ldots, 1)$, a quadratic function
	\[ \varphi_{\quadratic}(x) = x^\top \one - 4 - 1.25 (x(1) - x(2))^2 \]
	and a function stemming from applications in finance~\cite{bassamboo_portfolio_2008}:
	\[ \varphi_{\fin}(x) = \sum_{j=3}^d \Indicator{\phi(x(1),x(2),x(j)) \geq 0.5 \sqrt{d}} - 0.25 d - 0.1 \]
	where for any $(x, y, z) \in \R^3$, 
	\[\phi(x,y,x) = \left(0.25 \, x + 3(1 - 0.25^2)^{1/2} z \right) \sqrt{ F_{\Gamma(6,6)}^{-1}\left(F_{N}(y) \right)} \]
	with $F_{\Gamma(6,6)}$ the cumulative distribution function of the Gamma distribution with both shape and scale parameters set to $6$. The linear and quadratic test cases are standard functions considered in reliability analysis~\cite{ELMASRI2021, papaioannou_sequential_2016, PAPAKONSTANTINOU2023, uribe2021}. The financial function $\varphi_\fin$ is for instance considered in~\cite{bassamboo_portfolio_2008, chan_improved_2012, ELMASRI2021}: it brings additional insight compared to $\varphi_\lin$ and $\varphi_\quadratic$ because it is not differentiable. The two functions $\varphi_\lin$ and $\varphi_\quadratic$ have an intrinsic dimension equal to one and two, respectively, and thus satisfy Assumption~\ref{ass:FID}. Formally, $\varphi_\fin$ does not satisfy this assumption, but a closer look suggests that it does satisfy it approximately: indeed, if the components $X(i)$ are i.i.d., then the law of large numbers suggests that
	\[ \sum_{j=3}^d \Indicator{\phi(X(1),X(2),X(j)) \geq 0.5 \sqrt{d}} \approx d h(X(1), X(2)) \]
	 with $h(x,y) = \P(\phi(x,y, X) \geq 0.5 \sqrt d)$, suggesting that $\varphi_\fin$ approximately has an intrinsic dimension of (at most) two.
\end{description}


For each test function $\varphi$, different CE schemes lead to different estimators $\hat p$ of the form~\eqref{eq:IS}. Recall that $\nfin$ is the sample size used in this final estimation step, while $n$ refers to the sample size used to learn the auxiliary distribution in step~$2$ of Algorithms~\ref{alg:CE-proj} (CE with projection) and~\ref{alg:iCE-proj} (iCE with projection). To assess the accuracy of each estimator, we consider the distribution of the relative error $\lvert \hat p - p \rvert / p$. This unknown distribution is estimated by kernel regression (displayed as so-called violin plots) obtained by repeating each estimation $N = 200$ times. Table~\ref{tab:parameters} gathers the relevant parameters for the numerical simulations presented below.



\subsubsection{Numerical results}


Figures~\ref{fig:CE} and~\ref{fig:iCE} compare the performance of CE, CE-eig and CE-mean (Figure~\ref{fig:CE}) and of iCE, iCE-eig and iCE-mean (Figure~\ref{fig:iCE}). The goal on these two figures is to check the relevance of Theorem~\ref{thm:main} for true CE schemes and also to assess the influence of the choice of the projection directions. Let us first comment results for CE on Figure~\ref{fig:CE}.

% 
\begin{figure}[t]
    \centering
    \hspace*{0.1\linewidth}\includegraphics[width=0.5\linewidth]{legend_CE.png} \newline
    \begin{subfigure}[c]{.33\textwidth}
        \centering
        \includegraphics[width=\linewidth]{G_hyperplan_CE_p.png}
		%
        \includegraphics[width=\linewidth]{G_hyperplan_CE_spectrum.png}
        \caption{Results for $\varphi_\lin$.}
		\label{subfig:CE-lin}
    \end{subfigure}%
    \hfill
    \begin{subfigure}[c]{.33\textwidth}
        \centering
        \includegraphics[width=\linewidth]{G_quadratic_CE_p.png}
        \includegraphics[width=\linewidth]{G_quadratic_CE_spectrum.png}
        \caption{Results for $\varphi_\quadratic$.}
		\label{subfig:CE-quad}
    \end{subfigure}%
    \hfill
    \begin{subfigure}[c]{.33\textwidth}
        \centering
        \includegraphics[width=\linewidth]{G_portfolio_CE_p.png}
        \includegraphics[width=\linewidth]{G_portfolio_CE_spectrum.png}
        \caption{Results for $\varphi_\fin$.}
		\label{subfig:CE-fin}
    \end{subfigure}%
    \hfill
	\caption{Interpretation of the performance of CE, CE-eig and CE-mean via the spectral behavior. Results for CE on the test functions $\varphi_\quadratic$ and $\varphi_\fin$ are not displayed because this scheme did not converge. For similar reasons, results for CE is not displayed for $\lambda_{\max}(\hat \Sigma_4)$ on the test function $\varphi_\lin$ since most repetitions diverge beyond the third iteration. Top figures display the distribution of the relative error $\lvert \hat p - p \rvert / p$. Bottom figures show the distribution of $\lambda_{\min}(\hat \Sigma^\proj_t)$ and $\lambda_{\max}(\hat \Sigma_{t+1})$ for the first three iterations $t = 1,2$ and $3$.}
    \label{fig:CE}
\end{figure}
% 

The top row shows that CE-mean is the best algorithm, with the density of the relative error $\lvert \hat p - p \rvert / p$ more concentrated toward $0$. On the three test cases, CE and CE-eig yield very poor results: algorithms either do not converge, or yield a relative error close to $100\%$. In contrast, CE-mean gives much better results. Performance for $\varphi_\lin$ and $\varphi_\fin$ is satisfactory. The quadratic case $\varphi_\quadratic$ remains challenging, but there is a substantial improvement from CE-eig to CE-mean.

The second row of Figure~\ref{fig:CE} gives an interpretation of these results by looking at the spectrum of the different covariance matrices during the first three iterations of the CE scheme. These results corroborate the previous discussion in Section~\ref{sub:implications}: we see that the best-performing algorithm, CE-mean, is also the one with the largest smallest eigenvalue $\lambda_{\min}(\hat \Sigma^\proj_t)$. Moreover, these numerical results suggest that the simplified model~\eqref{eq:hat-Sigma} yields relevant insight into CE schemes: indeed, we see that a small $\lambda_{\min}(\hat \Sigma^\proj_t)$ leads to a large $\lambda_{\max}(\hat \Sigma_{t+1})$, which is exactly the behavior pointed out in Theorem~\ref{thm:main}. The idea is that when the smallest eigenvalue of the covariance matrix of the auxiliary distribution is small, the sample size required for accurate estimation becomes very large. Although, as mentioned previously, the projection step typically only acts on small variance terms and makes $\lambda_{\max}(\hat \Sigma^\proj_{t+1}) = 1$, the fact that $\lambda_{\max}(\hat \Sigma_{t+1})$ is very large suggests that $\hat \Sigma_{t+1}$ does not satisfactorily estimate $\Sigma_{t+1}$. It is interesting to observe how this estimation error then propagates in the successive iterations: indeed, if $\Sigma_{t+1}$ is poorly estimated, then one may expect that $\hat \Sigma^\proj_{t+1}$ will also be rigged with noise. But as $\hat \Sigma^\proj_{t+1}$ is the covariance matrix of the next auxiliary distribution $\hat g_{t+1}$, this error is expected to propagate in the next iteration. And indeed, we clearly see in Figure~\ref{fig:CE} that poorly-performing schemes such as CE-eig have vanishing $\lambda_{\min}(\hat \Sigma^\proj_t)$ and diverging $\lambda_{\max}(\hat \Sigma_{t+1})$. Note that small $\lambda_{\min}$'s may seem like a good idea since the optimal covariance matrix $\Sigma_A$ has $\lambda_{\min}(\Sigma_A) \approx 0.0075$. Thus, the poor behavior of CE-eig may seem surprising at first glance, but Theorem~\ref{thm:main} sheds light on this phenomenon: a small $\lambda_{\min}$ gets closer to the optimal value $\lambda_{\min}(\Sigma_A)$ but also entails a more difficult estimation problem, requiring large sample size.
\\

Figure~\ref{fig:iCE}
% 
\begin{figure}[!t]
    \centering
    \hspace*{0.1\linewidth}\includegraphics[width=0.5\linewidth]{legend_iCE.png} \newline
    \begin{subfigure}[c]{.33\textwidth}
        \centering
        \includegraphics[width=\linewidth]{G_hyperplan_iCE_p.png}
		%
        \includegraphics[width=\linewidth]{G_hyperplan_iCE_spectrum.png}
        \caption{Results for $\varphi_\lin$.}
		\label{subfig:iCE-lin}
    \end{subfigure}%
    \hfill
    \begin{subfigure}[c]{.33\textwidth}
        \centering
        \includegraphics[width=\linewidth]{G_quadratic_iCE_p.png}
        \includegraphics[width=\linewidth]{G_quadratic_iCE_spectrum.png}
        \caption{Results for $\varphi_\quadratic$.}
		\label{subfig:iCE-quad}
    \end{subfigure}%
    \hfill
    \begin{subfigure}[c]{.33\textwidth}
        \centering
        \includegraphics[width=\linewidth]{G_portfolio_iCE_p.png}
        \includegraphics[width=\linewidth]{G_portfolio_iCE_spectrum.png}
        \caption{Results for $\varphi_\fin$.}
		\label{subfig:iCE-fin}
    \end{subfigure}%
    \hfill
	\caption{Interpretation of the performance of iCE, iCE-eig and iCE-mean via the spectral behavior. Top figures display the distribution of the relative error $\lvert \hat p - p \rvert / p$. Bottom figures show the distribution of $\lambda_{\min}(\hat \Sigma^\proj_t)$ and $\lambda_{\max}(\hat \Sigma_{t+1})$ for the first three iterations $t = 1,2$ and~$3$.}
    \label{fig:iCE}
\end{figure}
% 
shows that this insight is also confirmed by the behavior of iCE. We see on the top row of this figure that iCE behaves poorly, thereby justifying the need for projection. When considering projections, we see that iCE-mean offers significant improvements over iCE-eig for the two test cases $\varphi_\quadratic$ and $\varphi_\fin$ (Figures~\ref{subfig:iCE-quad} and~\ref{subfig:iCE-fin}). For these two examples, the same interpretation as for CE holds: iCE-mean is the variant with the largest $\lambda_{\min}(\hat \Sigma^\proj_t)$, and during the iterations of the different schemes, we see that a decreasing $\lambda_{\min}(\hat \Sigma^\proj_t)$ translates into an increasing $\lambda_{\max}(\hat \Sigma_{t+1})$. Finally, the linear test case (Figure~\ref{subfig:iCE-lin}) shows that the insight is quite sharp. In this case, iCE-eig and iCE-mean have similar median results, but the distribution of the relative error is slightly more concentrated for iCE-eig, which thus displays favorable performance over iCE-mean. Albeit subtle, this better performance is visible on the spectrum behavior, with iCE-eig eventually leading to a larger $\lambda_{\min}(\hat \Sigma^\proj_t)$ than iCE-mean.







\section{Opening discussion}\label{sec:conclu}
 
 
Overall, the results presented in Figures~\ref{fig:CE} and~\ref{fig:iCE} tend to confirm the central insight from Theorem~\ref{thm:main}, namely that CE schemes with smaller $\lambda_{\min}(\Sigma)$ tend to perform better. We now discuss some potentially fruitful research directions stemming from this new way to look at CE schemes.

First, we have so far focused on the influence of projecting, but simply comparing CE and iCE without projection suggests that our results may reach further. The performance of CE compared to that of iCE is displayed in Figure~\ref{fig:CEvsiCE}
% 
\begin{figure}
	\centering
	\begin{subfigure}[c]{.45\textwidth}
		\includegraphics[width=\textwidth]{G_hyperplan_CE_iCE_p.png}
		\caption{Relative error.} \label{fig:CEvsiCE-a}
	\end{subfigure}
	\hfill
	\begin{subfigure}[c]{.45\textwidth}
		\includegraphics[width=\textwidth]{G_hyperplan_CE_iCE_spectrum.png}
		\caption{Evolution of the extreme eigenvalues.} \label{fig:CEvsiCE-b}
	\end{subfigure}
	\caption{Comparison of CE and iCE on the linear test function $\varphi_\lin$. (a) Distribution of the relative error $\lvert \hat p - p \rvert / p$. (b) Evolution of $\lambda_{\min}(\hat \Sigma^\proj_t)$ and $\lambda_{\max}(\hat \Sigma_{t+1})$ during the first three iterations. Most repetitions diverge beyond the third iteration for CE so that results of $\lambda_{\max}(\hat \Sigma_4)$ are not displayed.}
	\label{fig:CEvsiCE}
\end{figure}
% 
for the linear test case $\varphi_\lin$ (the only one on which CE without projection converges). We recover the well-known fact that iCE performs much better than CE (Figure~\ref{fig:CEvsiCE-a}), but it is striking to note that even in this case without projection, our insight continues to hold since iCE leads to a larger $\lambda_{\min}(\hat \Sigma^\proj_t)$ than CE (Figure~\ref{fig:CEvsiCE-b}). Again, this is surprising as iCE was not motivated by spectral considerations. Thus, it would be interesting to understand the reason why the smoothing of the indicator function (which is the primary feature of iCE) impacts the spectrum.

Second, although our results allow to understand the potential benefits of projection methods in CE schemes, they actually challenge the relevance of importance sampling in high dimension. Indeed, consider the following corollary to Theorem~\ref{thm:main}, directly obtained with $g = f$.
% 
\begin{corollary} \label{cor:monte-carlo}
	Consider the Monte Carlo scheme, i.e., $\Sigma = I$, in the regime $n = d^\kappa$ for some $\kappa > 0$. If Assumption~\ref{ass:FID} holds and $\inf_d p > 0$, then $\lVert \hat \Sigma_A - \Sigma_A \rVert_\op \Rightarrow 0$ for any $\kappa > 1$.
\end{corollary}
% 
\noindent Thus on the one hand, Theorem~\ref{thm:main} and Proposition~\ref{prop:range} suggest that importance sampling with auxiliary distribution $g = N(0, \Sigma)$ may require $n \gg d^{1/\lambda_{\min}(\Sigma)}$ to be consistent, even in good cases where $V \subset U$, while Corollary~\ref{cor:monte-carlo} suggests that $n \gg d$ is sufficient for Monte Carlo. Overall, this suggests that Monte Carlo behaves better than importance sampling in high dimension. Of course, this is contradicted by numerous numerical results which show that, in practice, Monte Carlo fails to estimate small probabilities in high dimension while importance sampling with a suitable auxiliary distribution may succeed. Thus, our results call for a theoretical justification of the improvement of importance sampling over Monte Carlo in high dimension which, in our view, is still lacking.


Third, our results shed light and complement earlier results. Consider in particular the following theoretical result on CE without projection.
% 
\begin{thm} [Theorem $2.2$ in~\cite{beh2023insight}] \label{thm:CE}
	Consider the notation introduced in Algorithm~\ref{alg:CE}, and assume that:
	\begin{itemize}
		\item $\inf_d p$ and $\inf_d \rho > 0$;
		\item $m, \nfin \to \infty$;
		\item for every $d$, $\varphi^{-1}(\{x\})$ has zero Lebesgue measure for every $x \in \R$.
	\end{itemize}
	Then for every $t \geq 0$, there exists $\kappa_t \in (0,\infty)$ such that if $\npar \gg d^{\kappa_t}$, then $\hat p / p \Rightarrow 1$ with $\hat p$ defined through~\eqref{eq:IS} with $g = \hat g_t$.
\end{thm}
% 
\noindent The difference with Theorem~\ref{thm:main} is that this result applies to the true CE scheme as per Algorithm~\ref{alg:CE}, whereas Theorem~\ref{thm:main} only addresses a related random matrix model. Its drawback however is that only existence of some exponent $\kappa_t$ is proved, but its value is not provided, and moreover, the result only shows the sufficiency of the condition $n \gg d^{\kappa_t}$. Furthermore, the authors argue in~\cite{beh2023insight} for a typical dependency of the form
\begin{align} \label{eq:kappa-conjecture}
	\kappa_t \propto \frac{1}{\displaystyle \min \left(\lambda_{\min}(\Sigma_1), \dots, \lambda_{\min}(\Sigma_{t-1}) \right)},
\end{align}
which is clearly supported by Theorem~\ref{thm:main}. It would interesting to see if the techniques developed here could be integrated into a global analysis of CE to give an explicit expression of $\kappa_t$ in Theorem~\ref{thm:CE} above, and maybe also show that the growth rate $d^{\kappa_t}$ is necessary.

Finally, although we believe that the spectral behavior of the estimated covariance matrices is an important ingredient for explaining the performance of a given iterative importance sampling scheme, we emphasize that this is not necessarily the whole story. For instance, Uribe et al.~\cite{uribe2021} have proposed projecting on the so-called Failure Informed Subspace (FIS). Preliminary results suggest that a naive implementation of this method as in Algorithm~\ref{alg:iCE-proj} confirms the insight gained previously but that a smarter numerical implementation leads to significant improvements over the algorithms considered here, while maintaining lower $\lambda_{\min}(\hat \Sigma^\proj_t)$. Several distinctive features of FIS may explain this: for instance the fact that, when projecting on FIS, the projection step $\Proj_r(\hat \Sigma_{t+1}, v)$ can be performed without having to compute the full covariance matrix $\hat \Sigma_{t+1}$; or the fact that FIS requires a differentiable function $\varphi$ and leverages additional information on its gradient. Investigating this behavior in more details appears to be a potentially fruitful research directions, both theoretically and numerically.



\section{Proof of Theorem~\ref{thm:main}} \label{sec:proofs}

In the rest of this section, we assume that Assumptions~\ref{ass:Sigma}, ~\ref{ass:FID} and~\ref{ass:MW} hold, that $\inf_d p > 0$, that either $V \subset U$ or $V \subset U_\perp$ and that $n = d^\kappa$ for some $\kappa > 0$. The proof is organized as follows. We first introduce in Section~\ref{sub:preliminary} some key objects and preliminary results that will be used throughout the proof. The proof is then decomposed in four steps:
\begin{description}
	\item[Step $1$ (Section~\ref{sub:step-1}):] we prove that $\lambda_{\max}(\hat \Sigma_A) \Rightarrow \infty$ when $\kappa < \kappa_*$, and that it implies that $\E \lVert \hat \Sigma_A - \Sigma_A \rVert_\op \to \infty$;
	\item[Step $2$ (Section~\ref{sub:step-2}):] we prove that $\lVert \hat \Sigma_A - \Sigma_A \rVert_\op \Rightarrow 0$ when $\kappa > \kappa_*$;
	\item[Step $3$ (Section~\ref{sub:step-3}):] we prove that $\kappa_* = 1/\lambda_1$ when $V \subset U_\perp$;
	\item[Step $4$ (Section~\ref{sub:step-4}):] we prove that $1 \leq \kappa_* \leq 1/\lambda_1$ when $V \subset U$.
\end{description}

Note that step 2 is the most delicate step which rests on recent results for the concentration of sums of independent random matrices~\cite{Brailovskaya24-0}.



\subsection{Preliminary results} \label{sub:preliminary}

Before proceeding to the proof we collect some results and introduce new objects that will be used in the sequel.

\begin{lemma}\label{lemma:ell}
	For any $x \in \R^d$ we have
	\begin{equation} \label{eq:expression-ell}
		\ell(x) = \lvert \Sigma \rvert^{1/2} = \lvert \Sigma \rvert^{1/2} \exp \left( \frac{1}{2} \sum_{k=1}^r \left( \frac{1}{\lambda_k} - 1 \right) \langle v_k, x \rangle^2 \right).
	\end{equation}
	In particular, $\ell(x) = \ell(P_V x)$.
\end{lemma}

\begin{proof}
	The expression~\eqref{eq:expression-ell} follows readily from the fact that $\ell = f/g$ and that
	\[ \Sigma^{-1} = \sum_{k=1}^r (\lambda^{-1}_k - 1) v_k v_k^\top + I, \]
	as a direct consequence of~\eqref{eq:Sigma}.
\end{proof}

In the sequel, it will often be useful to rewrite $\hat \Sigma_A$ in~\eqref{eq:Sigma} as
\begin{equation} \label{eq:Sigma-prime}
	\hat \Sigma_A = \frac{1}{np} \sum_{j=1}^{n'} \ell(X^\prime_i) X^\prime_i X^{\prime\top}_i - \mu_A \mu_A^\top
\end{equation}
where $n^\prime = \sum_{i=1}^n \xi_A(X_i)$ is a binomial random variable with the two parameters $n$ and $q = \P_g(X \in A)$, and conditionally on $n'$, the $X^\prime_i$'s are i.i.d.\ drawn according to~$g|_A$.

\begin{lemma}\label{lemma:q-n'}
	If $\inf_d p > 0$, then $\inf_d q > 0$. Moreover, $n'/(nq) \to 1$ in \textnormal{L}\textsubscript2.
\end{lemma}

\begin{proof}
	The arguments closely follow those of the proof of~\cite[Corollary $3.5$]{beh2023insight}, which we detail for completeness. We have
	\[ q = \P_g(X \in A) = \E_f \left( \frac{g(Y)}{f(Y)} \xi_A(Y) \right) = p \E_{f|_A} \left( \frac{g(Y)}{f(Y)} \right) = \E \left( 1/\ell_A(Y') \right) \]
	with $\ell_A = f|_A / g$ and $Y' \sim f|_A$, leading to $q = \E \left[ \exp \left( -\log \ell_A(Y') \right) \right]$. Introducing $D = \E(\log \ell_A(Y'))$ the Kullback--Leibler divergence between $f|_A$ and $g$ and using Jensen's inequality for the exponential function, we obtain $q \geq e^{-D}$. According to~\cite[Lemma~$3.4$]{beh2023insight}, we have
	\[ D = -\log p - \Psi(\Sigma_A) - \frac{1}{2} \lVert \mu_A \rVert^2 + \Psi(\Sigma^{-1} \Sigma_A) + \frac{1}{2} \mu_A^\top \Sigma^{-1} \mu_A \]
	with $\Psi(M) = (\Tr(M) - \log \lvert M \rvert - d)/2$ for any $d\times d$ matrix M. It is easily seen that $\Psi(M) \geq 0$ for any symmetric matrix $M$ (see for instance~\cite[Eq.~(10)]{beh2023insight}), so that
	\[ q \geq \exp\left( - \Psi(\Sigma^{-1} \Sigma_A) - \frac{1}{2} \mu_A^\top \Sigma^{-1} \mu_A \right). \]
	We have $\mu_A^\top \Sigma^{-1} \mu_A \leq \lVert \mu_A \rVert^2 / \lambda_1$ which is bounded ($\lambda_1 > 0$ is fixed, while $\sup_d \lVert \mu_A \rVert < \infty$ by~\eqref{eq:inf-p}). Moreover, $\Psi(\Sigma_A)$ is bounded by~\cite[Corollary~$4.2$]{beh2023insight} and so is $\Psi(\Sigma^{-1})$ by definition of $\Sigma$ and under Assumption~\ref{ass:Sigma}. Therefore, $\Psi(\Sigma^{-1} \Sigma_A)$ is bounded by~\cite[Lemma~$3.8$]{beh2023insight} and so $\inf_d q > 0$ as desired. Finally, the convergence $n'/(qn) \to 1$ in L\textsubscript{2} is immediate since $n'$ is a binomial random variable with parameter $(n,q)$.
\end{proof}

We will consider
\begin{equation} \label{eq:M}
	M = \frac{d}{n} \max_{1 \leq i \leq n} \xi_A(X_i) \ell(X_i) = \frac{1}{n^{1-1/\kappa}} \max_{1 \leq j \leq n'} \ell(X'_j)
\end{equation}
as well as
\begin{equation} \label{eq:W}
	W = \left\{ \begin{array}{ll} U & \text{if } V \subset U\\V & \text { if } V \subset U_\perp \end{array} \right.
\end{equation}
together with $(w_i)$ an orthonormal family spanning $W$. For future use we collect the following properties related to $W$.

\begin{lemma}\label{lemma:W}
	If Assumptions~\ref{ass:Sigma} and~\ref{ass:FID} hold and $V \subset U$ or $V \subset U_\perp$, then the following properties hold where $X' \sim g|_A$:
	\begin{itemize}
		\item $\sup_d \dim(W) < \infty$;
		\item $\ell(x) = \ell(P_W x)$ for any $x \in \R^d$;
		\item $P_W X^\prime$ and $P_{W_\perp} X^\prime$ are independent;
		\item if $V \subset U$, then $P_{W_\perp} X^\prime \sim N(0, P_{W_\perp})$;
		\item if $V \subset U_\perp$, then $P_{W_\perp} X^\prime \sim N(0, P_{W_\perp})|_A$.
	\end{itemize}
\end{lemma}

\begin{proof}
	The first property is obvious from the definition of $W$ and Assumptions~\ref{ass:Sigma} and~\ref{ass:FID}. For the second property, we note that $V \subset W$ so that $P_V P_W = P_V$: thus, Lemma~\ref{lemma:ell} gives $\ell(P_W x) = \ell(P_V P_W x) = \ell(P_V x) = \ell(x)$.
	
	Let us now consider $X' \sim g|_A$ and prove that $P_W X'$ and $P_{W_\perp} X'$ are independent. First, note that $P_W X$ and $P_{W_\perp} X$ are independent when $X \sim g$. Indeed, for $X \sim g$ we have $\Cov(P_W X, P_{W_\perp} X) = P_W \Sigma P_{W_\perp}$. Since $W_\perp \subset V_\perp$, it follows that $P_{W_\perp} v = P_{W_\perp} P_{V_\perp} v = 0$ for any $v \in V$. In view of the definition~~\eqref{eq:Sigma} of $\Sigma$, this implies that $\Sigma P_{W_\perp} = P_{W_\perp}$ and so $P_W \Sigma P_{W_\perp} = P_W P_{W_\perp} = 0$. This shows that $\Cov(P_W X, P_{W_\perp} X) = 0$ which implies the independence between $P_W X$ and $P_{W_\perp}$ as $X$ is Gaussian.
	
	Let us now prove that $P_W X^\prime$ and $P_{W_\perp} X^\prime$ are also independent. Let $\psi, \phi: \R^d \to \R_+$ be any measurable functions: by definition of $X^\prime$ we have
	\[ \E(\psi(P_W X^\prime) \phi(P_{W_\perp} X^\prime)) = \E_g(\psi(P_W X) \phi(P_{W_\perp} X) \mid X \in A) \]
	and so under Assumption~\ref{ass:FID} we obtain
	\[ \E(\psi(P_W X^\prime) \phi(P_{W_\perp} X^\prime)) = \left\{\begin{array}{ll} \E_g(\psi(P_U X) \phi(P_{U_\perp} X) \mid P_U X \in A) & \text{if } V \subset U,\\\E_g(\psi(P_V X) \phi(P_{V_\perp} X) \mid P_U X \in A) & \text{if } V \subset U_\perp. \end{array} \right. \]
	When $V \subset U_\perp$, we have $P_U X = P_U P_{V_\perp} X$ and so $P_U X \in A \Leftrightarrow P_{V_\perp} X \in A$ by~\eqref{eq:FID}. Therefore, in view of the previous display and the definition of $W$, we get
	\begin{equation} \label{eq:interm}
		\E(\psi(P_W X^\prime) \phi(P_{W_\perp} X^\prime)) = \left\{\begin{array}{ll} \E_g(\psi(P_W X) \phi(P_{W_\perp} X) \mid P_W X \in A) & \text{if } V \subset U,\\\E_g(\psi(P_W X) \phi(P_{W_\perp} X) \mid P_{W_\perp} X \in A) & \text{if } V \subset U_\perp, \end{array} \right.
	\end{equation}
	and so the independence between $P_W X'$ and $P_{W_\perp} X'$ follows from that between $P_W X$ and $P_{W_\perp} X$.
	
	Moreover, it follows from the previous expression that $P_{W_\perp} X'$ is equal in distribution to $P_{W_\perp} X$ when $V \subset U$ and to $P_{W_\perp} X \mid P_{W_\perp} X \in A$ when $V \subset U_\perp$. Therefore, in order to conclude the proof, it remains to prove that $P_{W_\perp} X \sim N(0, P_{W_\perp})$. Since $X$ is centered, we have to check that the variances agree. But since we have proved that $\Sigma P_{W_\perp} = P_{W_\perp}$, this readily follows: $\Var(P_{W_\perp} X) = P_{W_\perp} \Sigma P_{W_\perp} = P_{W_\perp}^2 = P_{W_\perp}$.
\end{proof}




\subsection{Step $1$: $\lambda_{\max}(\hat \Sigma_A) \Rightarrow +\infty$ when $\kappa < \kappa_*$} \label{sub:step-1}

First we assume that $\kappa < \kappa_*$ and we prove that $\lambda_{\max}(\hat \Sigma_A) \Rightarrow +\infty$. Note that this implies the result stated in Theorem~\ref{thm:main}, because $\lVert \hat \Sigma_A - \Sigma_A \rVert_\op \geq \lambda_{\max}(\hat \Sigma_A) - \lambda_{\max}(\Sigma_A)$ and $\sup_d \lambda_{\max}(\Sigma_A) < \infty$ by~\eqref{eq:inf-p}. Let
\[ J = \arg \max_{i=1, \ldots, n} \xi_A(X_i) \ell(X_i),\ x = \frac{X_J}{\lVert X_J \rVert} \ \text{ and } \ \chi = \frac{1}{d} \lVert P_{W_\perp} X_J \rVert^2. \]
Note that the probability of the event $\cap_{i=1}^n \{\xi_A(X_i) = 0\}$ vanishes because
\[ \P \left( \xi_A(X_i) = 0, i = 1, \ldots, n \right) = (1-q)^n \]
which goes to $0$ according to Lemma~\ref{lemma:q-n'}. Thus, without loss of generality, we assume that $J$ is well-defined. Moreover, this also implies that $\xi_A(X_J) = 1$. Starting from~\eqref{eq:hat-Sigma} we get
\begin{align*}
	x^\top \hat \Sigma_A x & = \frac{1}{np} \sum_{j=1}^n \xi_A(X_i) \ell(X_i) \langle X_i, x \rangle^2 - x^\top \mu_A \mu_A^\top x\\
	& \geq \frac{1}{n} \xi_A(X_J) \ell(X_J) \langle X_J, x \rangle^2 - \langle x, \mu_A \rangle^2\\
	& \geq \frac{1}{n} \ell(X_J) \lVert X_J \rVert^2 - \lVert \mu_A \rVert^2
\end{align*}
where the first inequality comes from the definition of $J$ and the second equality follows from the fact that $\xi_A(X_J) = 1$, $\langle X_J, x \rangle^2 = \lVert X_J \rVert^2$ and $\langle \mu_A, x \rangle^2 \leq \lVert \mu_A \rVert^2$. Recalling the definition~\eqref{eq:M} of $M$, we finally get
\[ \lambda_{\max}(\hat \Sigma_A) \geq x^\top \hat \Sigma_A x \geq M \chi - \lVert \mu_A \rVert^2. \]
Since $\sup_d \lVert \mu_A \rVert^2 < \infty$ by~\eqref{eq:inf-p}, it is enough to show that $M \chi \Rightarrow +\infty$. For any $\alpha > 0$ and $0 < \beta < 1$ we have
\[ \P \left( \chi M \leq \alpha \right) \leq \P \left( M \leq \alpha/\beta \right) + \P \left( \chi \leq \beta \right). \]
Since $\kappa < \kappa_*$, we have
\[ 1 - \frac{1}{\kappa} < 1 - \frac{1}{\kappa_*} = \gamma_* \]
and so according to Assumption~\ref{ass:MW}, we have $M \Rightarrow +\infty$ which implies that $\P \left( M \leq 2\alpha \right) \to 0$. To control $\lVert P_{W_\perp} X_J \rVert$, let us use~\eqref{eq:Sigma-prime} to rewrite $X_J$ as
\[ X_J = \arg \max_{1 \leq j \leq n^\prime} \ell(X^\prime_j). \]
If $\P^\prime = \P(\ \cdot \mid n^\prime)$, then by exchangeability we get
\begin{align*}
	\P^\prime \left( \chi \leq \beta \right) & = \P^\prime \left( \frac{1}{d} \lVert P_{W_\perp} X^\prime_1 \rVert^2 \leq \beta \mid \ell(X^\prime_1) \geq \ell(X^\prime_j), j = 2, \ldots, n^\prime \right)\\
	& = \P^\prime \left( \frac{1}{d} \lVert P_{W_\perp} X^\prime_1 \rVert^2 \leq \beta \mid \ell(P_W  X^\prime_1) \geq \ell(X^\prime_j), j = 2, \ldots, n^\prime \right)
\end{align*}
using $\ell(x) = \ell(P_W x)$ by Lemma~\ref{lemma:W}. Thus, using the independence between $P_{W_\perp} X^\prime_1$, $P_W X^\prime_1$ (by Lemma~\ref{lemma:W}) and the other $X^\prime_j$'s, we obtain
\[ \P^\prime \left( \chi \leq \beta \right) = \P^\prime \left( \frac{1}{d} \lVert P_{W_\perp} X^\prime_1 \rVert^2 \leq \beta \right) = \P \left( \frac{1}{d} \lVert P_{W_\perp} X^\prime_1 \rVert^2 \leq \beta \right) \]
with the last equality coming from the independence between $n^\prime$ and $X^\prime_1$. Integrating over $n^\prime$ and recalling that $X^\prime_1 \sim g|_A$ we thus obtain
\begin{align*}
	\P \left( \chi \leq \beta \right) & = \P_g \left( \frac{1}{d} \lVert P_{W_\perp} X \rVert^2 \leq \beta \mid X \in A \right)\\
	& \leq \frac{1}{q} \P_g \left( \frac{1}{d} \lVert P_{W_\perp} X \rVert^2 \leq \beta \right)\\
	& = \frac{1}{q} \P_f \left( \frac{1}{d} \lVert P_{W_\perp} \Sigma^{1/2} Y \rVert^2 \leq \beta \right)
\end{align*}
using for the last equality that $\Sigma^{1/2} Y \sim g$ if $Y \sim f$. Since $P_{W_\perp} \Sigma^{1/2} = P_{W_\perp}$ (which comes from $W_\perp \subset V_\perp$ along the same arguments as those leading to $\Sigma P_{W_\perp} = P_{W_\perp}$ in the proof of Lemma~\ref{lemma:W}), we finally obtain
\[ \P \left( \chi \leq \beta \right) \leq \frac{1}{q} \P_f \left( \frac{1}{d} \lVert P_{W_\perp} Y \rVert^2 \leq \beta \right). \]
Since $\lVert P_{W_\perp} Y \rVert^2$ follows a chi-square distribution with $d - \dim(W)$ degrees of freedom and $\sup_d \dim(W)< \infty$, we get $\frac{1}{d} \lVert P_{W_\perp} Y \rVert^2 \Rightarrow 1$ and so the probability in the right-hand side of the previous display vanishes since $\beta < 1$. Since finally $\inf_d q > 0$ by Lemma~\ref{lemma:q-n'} we get that $\P(\chi \leq \beta) \to 0$ which achieves the proof of this step.







\subsection{Step $2$: $\lVert \hat \Sigma_A - \Sigma_A \rVert_\op \Rightarrow 0$ when $\kappa > \kappa_*$} \label{sub:step-2}

Now we assume that $\kappa > \kappa_*$ and we prove that $\lVert \hat \Sigma_A - \Sigma_A \rVert_\op \Rightarrow 0$. Let $x \in \R^d$ with $\lVert x \rVert = 1$ and $\Delta = \hat \Sigma_A - \Sigma_A$. Then writing $x = y + y_\perp$ with $y \in W$ and $y_\perp \in W_\perp$ we obtain
\[ x^\top (\hat \Sigma_A - \Sigma_A) x = y^\top \Delta y + y^\top_\perp \Delta y_\perp + y^\top_\perp \Delta y + y^\top \Delta y_\perp \]
and so, since $\lVert y \rVert, \lVert y_\perp \rVert \leq 1$,
\[ \sup_{x \in \R^d: \lVert x \rVert = 1} \lvert x^\top (\hat \Sigma_A - \Sigma_A) x \rvert \leq \sup_{y \in W^*} \left \lvert y^\top \Delta y \right \rvert + \sup_{y_\perp \in W^*_\perp} \left \lvert y_\perp^\top \Delta y_\perp \right \rvert + 2\sup_{\substack{y \in W^*\\y_\perp \in W^*_\perp}} \left \lvert y^\top \Delta y_\perp \right \rvert \]
with $W^* = \{x \in W: \lVert x \rVert \leq 1\}$ and $W^*_\perp = \{x \in W_\perp: \lVert x \rVert \leq 1\}$. The following lemma plays an important role in the proofs.

\begin{lemma}\label{lemma:moments}
	Assume that Assumption~\ref{ass:Sigma} holds and let $X' \sim g|_A$. Then for any $w \in W^*$, any $\alpha < (1-\lambda_1)^{-1}$ and any $\beta > 0$, we have
	\[ \sup_d \E \left( \ell(X')^\alpha \lvert \langle w, X' \rangle \rvert^\beta \right) < \infty. \]
\end{lemma}

\begin{proof}
	Since $\E \left( \ell(X')^\alpha \lvert \langle w, X' \rangle \rvert^\beta \right) \leq q^{-1} \E_g \left( \ell(X)^\alpha \lvert \langle w, X \rangle \rvert^\beta \right)$ and $\inf_d q > 0$ by Lemma~\ref{lemma:q-n'}, it is enough by H\"older's inequality to prove that
	\[ \sup_d \E_g \left( \ell(X)^\alpha \right) < \infty \ \text{ and } \ \sup_d \E_g \left( \lvert \langle w, X \rangle \rvert^\beta \right) < \infty. \]
	Since $\ell = f/g$ we have $\E_g(\ell(X)^\alpha) = \E_f(\ell(Y)^{\alpha-1})$. Starting from~\eqref{eq:expression-ell}, we see that
	\[ \E_f(\ell(Y)^{\alpha-1}) = \lvert \Sigma \rvert^{(\alpha-1)/2} \E_f \left[ \exp \left( \frac{\alpha - 1}{2} \sum_{k=1}^r (\lambda_k^{-1}-1) \langle Y, v_k \rangle^2 \right) \right]. \]
	Since $Y$ is standard Gaussian and the $v_k$'s are orthonormal, the $Y^\top v_k$ are i.i.d.\ standard Gaussian random variables in dimension one, so that
	\[ \E_f(\ell(Y)^{\alpha-1}) = \lvert \Sigma \rvert^{(\alpha-1)/2} \prod_{k=1}^r \E \left(e^{ \frac{1}{2} (\alpha - 1)(\lambda_k^{-1} - 1) N^2} \right) \ \text{ where } \ N \sim N(0,1). \]
	Assumption~\ref{ass:Sigma} implies that $\sup_d \lvert \Sigma \rvert^{1/2} < \infty$, and from the density of the standard Gaussian random variable it follows immediately that
	\[ \E \left(e^{\frac{1}{2} (\alpha-1)(\lambda_k^{-1} - 1) N^2} \right) < \infty \Leftrightarrow (\alpha - 1)(\lambda_k^{-1} - 1) < 1 \]
	and so
	\[ \E_f(\ell(Y)^{\alpha-1}) < \infty \Leftrightarrow \max_k (\alpha - 1)(\lambda_k^{-1} - 1) = (\alpha - 1) (\lambda_1^{-1} - 1) < 1. \]
	This achieves to prove that $\E_g(\ell(X)^\alpha) < \infty$ for $\alpha < (1-\lambda_1)^{-1}$. Let us now show that $w^T X$ has finite moments of all orders, uniformly in $d$. By definition, $w^\top X$ is a centered Gaussian random variable with variance
	\[ \Var_g(w^\top X) = w^\top \Sigma w = \sum_{k=1}^r (\lambda_k - 1) \langle v_k, w \rangle^2 + \lVert w \rVert^2 \leq r \lambda_{\max}(\Sigma) + 1, \]
	using that $\lVert v_k \rVert = \lVert w \rVert = 1$. By Assumption~\ref{ass:Sigma} we therefore get $\sup_d \Var_g(w^\top X) < \infty$ which readily implies that all moments of $w^\top X$ are finite, uniformly in $d$.
\end{proof}


\subsubsection{Control of $y \in W^*$} 

Here we show that the term $\sup_{y \in W^*} \left \lvert y^\top \Delta y \right \rvert$ converges to $0$ in L\textsubscript1. Let $y \in W^*$: decomposing $y \in W$ on the orthornormal basis $(w_i)$ of $W$, we get
\[ y^\top \Delta y = \sum_{k, \ell = 1}^{\dim(W)} \langle y, w_k \rangle \langle y, w_\ell \rangle w_k^\top \Delta w_\ell \]
so that $\sup_{y \in W^*} \left \lvert y^\top \Delta y \right \rvert \leq \sum_{k, \ell} \lvert w_k^\top \Delta w_\ell \rvert$. In particular, since $\sup_d \dim(W) < \infty$, in order to prove the result it is enough to prove that $\lvert w_k^\top \Delta w_\ell \rvert \lone 0$ for any $1 \leq k, \ell \leq \dim(W)$. Since $\Delta = \hat \Sigma_A - \E(\hat \Sigma_A)$, we can write
\[ w_k^\top \Delta w_\ell = \frac{1}{np} \sum_{j=1}^{n'} \left( \chi_j - \E(\chi_j) \right) \ \text{ with } \ \chi_j = \ell(X'_j) \langle X'_j, w_k \rangle \langle X'_j, w_\ell \rangle. \]
A straightforward extension of Lemma~\ref{lemma:moments} implies that $\chi_1$ (and therefore $\chi_1 - \E(\chi_1)$) has finite moments of order $\eta$ for some fixed $\eta > 1$, uniformly in $d$. Recalling that $n'/(qn) \to 1$ in L\textsubscript2 and $\inf_d q > 0$ by Lemma~\ref{lemma:q-n'}, the next lemma implies that $w_k^\top \Delta w_\ell \lone 0$. For lack of reference to this probably standard result (a triangular version of the law of large numbers), we provide its proof.

\begin{lemma}\label{lem:triangular-lln}
	For each $d \geq 1$, let $\xi_{n,d} = \sum_{i=1}^n \zeta_{i,d}$ where $\zeta_{1,d}, \ldots, \zeta_{n,d}$ are centered and i.i.d.\ with $\sup_d \E\lvert \zeta_{1,d}\rvert^\eta < \infty$ for some $\eta > 1$. Then $n^{-1}\xi_{n,d} \Lone 0$.
\end{lemma}

\begin{proof}
	Let $\zeta_d = \zeta_{1,d}$ and $C = \sup_n \E \lvert \zeta_d \rvert^\eta$, which by assumption is finite. If $\lVert X \rVert_\eta = [\E(\lvert X \rvert^\eta)]^{1/\eta}$ denotes the $\eta$-norm, we have
	\[ \lVert \xi_{n,d} \rVert_\eta = \frac{1}{n} \Big \lVert \sum_{i=1}^n \zeta_{i,d} \Big \rVert_\eta \leq \lVert \zeta_d \rVert_\eta = [\E (\lvert \zeta_d \rvert^\eta)]^{1/\eta} \leq C^{1/\eta} \]
	which shows that $\xi_{n,d}$ is uniformly integrable. Thus, in order to show that $\xi_{n,d} \lone 0$, it is enough to show that $\xi_{n,d} \Rightarrow 0$. If $\eta \geq 2$ this is immediate by considering the variance, so let us assume $\eta < 2$. According to~\cite[Theorem 2.2.11]{Durrett_2019}, since the $\zeta_{i,d}$ are i.i.d., it is sufficient that
	\begin{itemize}
	\item[(i)] $n \P(\lvert \zeta_d \rvert > n) \to 0$,
	\item[(ii)] $\E[ \zeta_d ; \lvert \zeta_d \rvert \leq n] \to 0$, and
	\item[(iii)] $n^{-1} \E[ \zeta_d^2 ; \lvert \zeta_d \rvert \leq n] \to 0$.
	\end{itemize}
	The first two terms are bounded by $C/n^{\eta-1}$. Indeed, for (i) this directly comes from Markov inequality. For (ii), since $\E(\zeta_n) = 0$ it follows that $\E[ \zeta_d ; \lvert \zeta_d \rvert \leq n] = -\E[ \zeta_d ; \lvert \zeta_d \rvert \geq n]$ and so
	\[ \left \lvert \E[ \zeta_d ; \lvert \zeta_d \rvert \leq n]\right \rvert \leq \E[ \lvert \zeta_d \rvert ; \lvert \zeta_d \rvert \geq n] \leq \frac{C}{n^{\eta-1}}. \]
	For (iii), writing $\zeta_d^2 = \lvert \zeta_d \rvert^\eta \lvert \zeta_d \rvert^{2-\eta}$ and using $\lvert \zeta_d \rvert^{2-\eta} \leq n^{2-\eta}$ when $\lvert \zeta_d \rvert \leq n$ (and $\eta < 2$) leads to the bound $C / n^{\eta-1}$. In either case, we see that all three terms vanish as desired.
\end{proof}





\subsubsection{Control of $y \in W^*, y_\perp \in W^*_\perp$}

Let us now control the term $\sup_{y \in W^*, y_\perp \in W^*_\perp} \left \lvert y^\top \Delta y_\perp \right \rvert$. We begin with the following lemma.

\begin{lemma} \label{lemma:Sigma-A-mu-A-W}
	If Assumption~\ref{ass:FID} holds, then $y_\perp^\top \Sigma_A y = 0$ for any $y \in W$ and $y_\perp \in W_\perp$. Moreover, if $V \subset U$ then $P_{W_\perp} \mu_A = 0$ while if $V \subset U_\perp$ then $P_W \mu_A = 0$.
\end{lemma}

\begin{proof}
	First, note that $\mu_A \in U$, i.e., $P_{U_\perp} \mu_A = 0$: indeed,
	\[ P_{U_\perp} \mu_A = \E_f(P_{U_\perp} Y \mid Y \in A) = \E_f(P_{U_\perp} Y \mid P_U Y \in A) = \E_f(P_{U_\perp} Y) = 0, \]
	with the second equality coming from Assumption~\ref{ass:FID} and the third equality coming from the fact that $P_U Y$ and $P_{U_\perp} Y$ are independent (because $Y$ is standard Gaussian).
	
	Let us now prove that $y_\perp^\top \Sigma_A y = 0$ for $y \in W$ and $y_\perp \in W_\perp$. First consider the case where $W = U$. Consider $x_\perp \in U_\perp$, $x \in U$: then
	\begin{align*}
		x_\perp^\top \Sigma_A x & = \E_f ( \langle x, Y - \mu_A \rangle \langle x_\perp, Y - \mu_A \rangle \mid Y \in A)\\
		& = \E_f ( \langle x, P_U (Y - \mu_A) \rangle \langle x_\perp, P_{U_\perp} (Y - \mu_A) \rangle \mid P_U Y \in A)\\
		& = \E_f ( \langle x, P_U (Y - \mu_A) \rangle \mid P_U Y \in A) \times \E_f ( \langle x_\perp, P_{U_\perp} (Y - \mu_A) \rangle)
	\end{align*}
	which is $0$ since $\E_f ( \langle x, P_U (Y - \mu_A) \rangle \mid P_U Y \in A) = 0$ by linearity and definition of $\mu_A$. Note that the relation $x^\top_\perp \Sigma_A x = 0$ for any $x \in U$, $x_\perp \in U_\perp$ shows that $\Sigma_A$ leaves $U$ and $U_\perp$ invariant, i.e., $\Sigma_A x \in U$ (resp.\ $U_\perp$) if $x \in U$ (resp.\ $U_\perp$). Actually, a stronger result holds, namely $\Sigma_A x = x$ if $x \in U_\perp$. To prove this, it is enough to prove that $x^\top \Sigma_A x = \lVert x \rVert$ for $x \in U_\perp$. Let $x \in U_\perp$: then
	\[ x^\top \Sigma_A x = \E_f \left( \langle x, Y - \mu_A \rangle^2 \mid Y \in A \right) = \E_f \left( \langle x, P_{U_\perp} Y \rangle^2 \mid P_U Y \in A \right) \]
	using for the second equality that $P_{U_\perp} \mu_A = 0$ (and Assumption~\ref{ass:FID}). Since $Y$ is standard Gaussian, $P_U Y$ and $P_{U_\perp} Y$ are independent, leading to
	\[ x^\top \Sigma_A x = \E_f \left( \langle x, P_{U_\perp} Y \rangle^2 \right) = \lVert x \rVert^2 \]
	as desired.
	
	Let us now prove that $y_\perp^\top \Sigma_A y = 0$ for $y \in W$ and $y_\perp \in W_\perp$ in the case $W = V$ when $V \subset U_\perp$. Let $y \in V$ and $y_\perp \in V_\perp$: then $y \in U_\perp$ and so $\Sigma_A y = y$, so that $y^\top_\perp \Sigma_A y = y^\top_\perp y = 0$ as desired.
	
	To conclude the proof let us prove that either $P_W \mu_A = 0$ or $P_{W_\perp} \mu_A = 0$ depending on the case considered. Since we have proved $P_{U_\perp} \mu_A = 0$ this shows that $P_{W_\perp} \mu_A = 0$ when $W = U$. Let us finally consider $W = V$: since this is the definition of $W$ when $V \subset U_\perp$, we have $P_W \mu_A = P_V \mu_A = P_V P_{U_\perp} \mu_A = 0$. This concludes the proof.
\end{proof}

Fix in rest of this step $y \in W^*$, $y_\perp \in W^*_\perp$. Since $\Delta = \hat \Sigma_A - \Sigma_A$, Lemma~\ref{lemma:Sigma-A-mu-A-W} implies that $y^\top \Delta y_\perp = y^\top \hat \Sigma_A y_\perp$. Decomposing $y \in W$ on the orthornormal basis $(w_i)$ of $W$, we get
\[ y^\top \hat \Sigma_A y_\perp = \sum_{k=1}^{\dim(W)} \langle y, w_k \rangle w_k^\top \hat \Sigma_A y_\perp = \sum_{k=1}^{\dim(W)} \langle y, w_k \rangle \langle y_\perp, P_{W_\perp} \hat \Sigma_A w_k \rangle \]
and so since $\lVert y \rVert, \lVert y_\perp \rVert, \lVert w_k \rVert \leq 1$, we obtain
\[ \left \lvert y^\top \hat \Sigma_A y_\perp \right \rvert \leq \sum_{k=1}^{\dim(W)} \lVert P_{W_\perp} \hat \Sigma_A w_k \rVert. \]
Since $\sup_d \dim(W) < \infty$ it is therefore enough to show that $\lVert P_\perp \hat \Sigma_A w \rVert \Rightarrow 0$ for every $w \in W^*$. Starting from the expression~\eqref{eq:Sigma-prime} of $\hat \Sigma_A$, we have
\[ P_{W_\perp} \hat \Sigma_A w = \frac{1}{np} \sum_{j=1}^{n'} \ell(X'_j) \langle w, X'_j \rangle P_{W_\perp} X'_j - \langle \mu_A, w \rangle P_{W_\perp} \mu_A. \]
According to Lemma~\ref{lemma:Sigma-A-mu-A-W}, the second term in the right-hand side of the previous display vanishes (either $P_{W_\perp} \mu_A = 0$ or $\langle w, \mu_A \rangle = \langle w, P_W \mu_A \rangle = 0$), so that if we define $\ell'_j = \ell(X'_j) w^\top X'_j$ and $Z_j = P_{W_\perp} X'_j$, we can write
\[ P_{W_\perp} \hat \Sigma_A w = \frac{1}{np} \sum_{j=1}^{n'} \ell'_j Z_j. \]
Let in the rest of the proof $\tilde \P = \P(\ \cdot \ |\ n', P_W X'_j, j = 1, \ldots, n')$. We have
\[ \tilde \E \left( \lVert P_{W_\perp} \hat \Sigma_A w \rVert^2 \right) = \frac{1}{n^2p^2} \sum_{1 \leq i, j \leq n'} \ell'_i \ell'_j \E (Z_i^\top Z_j) \]
which we write $\tilde \E ( \lVert P_{W_\perp} \hat \Sigma_A w \rVert^2 ) = p^{-2}(D + CP)$ distinguishing between the diagonal and the cross-product terms:
\[ D = \frac{\E(\lVert Z_1 \rVert^2)}{n^2} \sum_{j = 1}^{n'} \ell^{\prime 2}_j \ \text{ and } \ CP = \frac{\lVert \E(Z_1) \rVert^2}{n^2} \sum_{1 \leq i \neq j \leq n'} \ell'_i \ell'_j. \]
First, note that $D \Rightarrow 0$, which we state as a separate lemma.

\begin{lemma}
	$D \Rightarrow 0$.
\end{lemma}

\begin{proof}
	Since $\lVert Z_1 \rVert = \lVert P_{W_\perp} X'_1 \rVert \leq \lVert X'_1 \rVert$ and $X'_1 \sim g|_A$, we get
	\[ \E\lVert Z_1 \rVert^2 \leq \frac{1}{q} \E_g \lVert X \rVert^2 = \frac{1}{q} \E_f \lVert \Sigma^{1/2} Y \rVert^2 \leq \frac{\lambda_{\max}(\Sigma)}{q} \E_f \lVert Y \rVert^2 = \frac{\lambda_{\max}(\Sigma)}{q} d.  \]
	Writing
	\[ \frac{1}{n} \sum_{j = 1}^{n'} \ell_j^{\prime 2} \leq \left( \max_{1 \leq j \leq n'} \ell(X'_j) \right) \bar \ell' \ \text{ with } \ \bar \ell' = \frac{1}{n} \sum_{j = 1}^{n'} \ell(X'_j) \langle w, X'_j \rangle^2 \]
	and recalling the definition~\eqref{eq:M} of $M$, we therefore get
	\[ D \leq q^{-1} \lambda_{\max}(\Sigma_A) M \bar \ell. \]
	We have $\inf_d q > 0$ by Lemma~\ref{lemma:q-n'} and $\sup_d \lambda_{\max}(\Sigma_A) < \infty$ by~\eqref{eq:inf-p}. Moreover, since $\kappa > \kappa_*$, we have $M \Rightarrow 0$ by Assumption~\ref{ass:MW}. Moreover, it readily follows from Lemma~\ref{lemma:moments} that $\bar \ell'$ is tight, so that the upper bound of the previous display vanishes, which implies $D \Rightarrow 0$ as desired.
\end{proof}

We now proceed to controlling the cross-product terms $CP$. The case $V \subset U$ is easy since in this case $CP = 0$ because $\E(Z_1) = \E(P_{W_\perp} X') = 0$ according to Lemma~\ref{lemma:W}. So consider now the case $V \subset U_\perp$. In this case $CP$ remains centered: indeed,
\[ \E(\ell'_1) = \E(\ell(X'_1) w^\top X'_1) = \E_g( \ell(X) w^\top X \mid X \in A) = \frac{1}{q} \E_g( \ell(X) w^\top X \xi_A(X)) \]
and recalling that $\ell = f/g$, we get
\[ \E(\ell(X'_1) w^\top X'_1) = \frac{1}{q} \E_f( w^\top Y \xi_A(Y)) = \frac{p}{q} w^\top \mu_A. \]
Since $P_W \mu_A = 0$ when $V \subset U_\perp$ by Lemma~\ref{lemma:Sigma-A-mu-A-W}, this shows that $\ell'_1$ and thus $CP$ is centered. Moreover, we have according to Lemma~\ref{lemma:W}
\[ \E(Z_1) = \E(P_{W_\perp} X') = \E_f(P_{W_\perp} Y \mid P_{W_\perp} Y \in A) = P_{V_\perp} \E_f(Y \mid P_{V_\perp} Y \in A) \]
using that $W = V$ when $V \subset U_\perp$. Then Assumption~\ref{ass:FID} gives
\[ \E(Z_1) = P_{V_\perp} \E_f(Y \mid P_U P_{V_\perp} Y \in A) \]
and since $U \subset V_\perp$ we finally obtain
\[ \E(Z_1) = P_{V_\perp} \E_f(Y \mid P_U Y \in A) = P_{V_\perp} \mu_A \]
and so $\lVert \E(Z_1) \rVert \leq \lVert \mu_A \rVert$. Since $\sup_d \lVert \mu_A \rVert < \infty$ by~\eqref{eq:inf-p}, we obtain $\sup_d \lVert \E(Z_1) \rVert < \infty$. Resuming the proof of $CP \Rightarrow 0$, we have $D + CP \geq 0$ so $D \geq -CP$ and since $D \geq 0$, this gives $D \geq (-CP)^+$. Moreover, Lemma~\ref{lemma:moments} implies that $\sup_d \E(\ell^{\prime \eta}_1) < \infty$ for some $\eta > 1$: by Minkowski's inequality, this implies that $\sup_d \E(\lvert CP \rvert^\eta) < \infty$ and thus that $CP$ is uniformly integrable. Gathering the above, we see that:
\begin{itemize}
	\item $D \geq (-CP)^+$;
	\item $CP$ is uniformly integrable;
	\item $\E(CP) = 0$;
	\item $D \Rightarrow 0$.
\end{itemize}
Therefore, $D$ and $-CP$ satisfy the assumptions of the next lemma, which achieves to prove that $CP \Rightarrow 0$.

\begin{lemma}
	Consider random variables $X,Y$ with $Y \geq X^+$, $X$ uniformly integrable, $\E(X) = 0$ and $Y \Rightarrow 0$: then $X \Lone 0$.
\end{lemma}

\begin{proof}
	$Y \geq X^+$ and $Y \Rightarrow 0$ implies $X^+ \Rightarrow 0$. Since $X$ is uniformly integrable, so is $X^+$ because $\lvert X \rvert \geq X^+$ and so $\E(X^+) \to 0$. Since $\E(X) = 0$, we have $\E(\lvert X \rvert) = 2\E(X^+)$ and so $X \to 0$ in L\textsubscript1.
\end{proof}

Since $D, CP \Rightarrow 0$ and $\inf_d p > 0$ we obtain going back to the relation $\tilde \E ( \lVert P_{W_\perp} \hat \Sigma_A w \rVert^2 ) = p^{-2}(D + CP)$ that $\tilde \E(\lVert P_{W_\perp} \hat \Sigma_A w \rVert^2) \Rightarrow 0$, which implies by dominated convergence that $\lVert P_{W_\perp} \hat \Sigma_A w \rVert \Rightarrow 0$ since for any $\varepsilon > 0$,
\[ \P \left( \lVert P_{W_\perp} \hat \Sigma_A w \rVert \geq \varepsilon \right) \leq \E \left[ \min \left( 1, \frac{1}{\varepsilon^2} \tilde \E(\lVert P_{W_\perp} \hat \Sigma_A w \rVert^2) \right) \right]. \]
This achieves to prove that $\sup_{y \in W^*, y_\perp \in W^*_\perp} \left \lvert y^\top \Delta y_\perp \right \rvert \Rightarrow 0$.


\subsubsection{Control of $y \in W^*_\perp$} \label{sec:y-V-perp}
Let us finally control the term $\sup_{y_\perp \in W^*_\perp} \left \lvert y_\perp^\top \Delta y_\perp \right \rvert$. Define
\[ S = \frac{1}{np} \sum_{j=1}^{n'} \ell(X'_j) (P_{W_\perp} X'_j) (P_{W_\perp} X^\prime_j)^\top \]
so that for $y_\perp \in W^*_\perp$, we have $y^\top_\perp \Delta y_\perp = y^\top_\perp (S' - \E(S')) y_\perp$. Further, we rewrite
\[ S = \sum_{j=1}^{n'} T_j T_j^\top \ \text{ with } \ \alpha_j = \frac{1}{np} \ell(X'_j) \ \text{ and } \ T_j = \sqrt{\alpha_j} P_{W_\perp} X'_j. \]
Recall that $\tilde \P = \P(\cdot \mid n', P_W X_j', j = 1, \ldots, n')$: a direct consequence of Lemma~\ref{lemma:W} is that under $\tilde \P$, the $T_j$'s are independent and $T_j \sim N(0, \Sigma_j) |_{B_j}$ where
\[ \Sigma_j = \alpha_j P_{W_\perp} \ \text{ and } \ B_j = \left\{ \begin{array}{ll}
	\R^d & \text{if } V \subset U,\\
	\sqrt{\alpha_j} A = \{\sqrt{\alpha_j} x: x \in A\} & \text{if } V \subset U_\perp.
\end{array} \right. \]
In the case $V \subset U$, we can therefore invoke~\cite[Theorem $3$.$21$]{Brailovskaya24-0} to get a bound on $\tilde \E \lVert S' - \E(S') \rVert_\op$. However, it is not difficult to adapt the proof of this result to get a bound in the case $V \subset U_\perp$ as well. For completeness, the generalization is established in the Appendix~\ref{appendix}.


\begin{thm} [{\protect{Extension of~\cite[Theorem~$3$.$21$]{Brailovskaya24-0}}}] \label{thm:B-vH-unconditioned}
	With the notation above, there exists a finite, universal constant $C > 0$ such that
    \begin{multline} \label{eq:bound-RMT}
        \tilde \E \lVert S - \tilde \E S \rVert_\op \leq \beta C \\
		\times \left( \Big\lVert \sum_{j=1}^{n'} \Tr(\Sigma_j) \Sigma_j \Big \rVert_\op^{1/2} + \left(\Big \lVert \sum_{j=1}^{n'} \Sigma_j^2 \Big \rVert_\op^{1/2} + \max_{j\leq n'}\Tr (\Sigma_j) \right)\log^3(d+n')\right)
    \end{multline}
	where $\beta = 1$ if $V \subset U$ and $\beta = p^{-1}$ if $V \subset U_\perp$.
\end{thm}

It remains to control the upper bound and show that it vanishes when $n = d^\kappa$ with $\kappa > \kappa_*$. By definition, $\Tr(\Sigma_j) = \alpha_j \Tr(P_{W_\perp}) = \alpha_j (d - \dim(W))$, which gives
\[ \sum_{j=1}^{n'} \Tr(\Sigma_j) \Sigma_j = \sum_{j=1}^{n'} (d - \dim(W)) \alpha^2_j P_{W_\perp} \]
and so
\[ \Big \lVert \sum_{j=1}^{n'} \Tr(\Sigma_j) \Sigma_j \Big \rVert_\op = \frac{d - \dim(W)}{n^2p^2} \sum_{j=1}^{n'} \ell(X'_j)^2 \lVert P_{W_\perp} \rVert_\op^{1/2} \leq \frac{d}{n^2p^2} \sum_{j=1}^{n'} \ell(X'_j)^2 \]
where we have used $\lVert P_{W_\perp} \rVert_\op = 1$. Writing
\[ \frac{d}{n^2} \sum_{j=1}^{n'} \ell(X'_j)^2 \leq \frac{d}{n} \max_{1 \leq j \leq n'} \ell(X'_j) \times \frac{1}{n} \sum_{j=1}^{n'} \ell(X'_j) \]
and recalling the definition~\eqref{eq:M} of $M$, we finally get
\[ \Big \lVert\sum_{j=1}^{n'} \Tr(\Sigma_j) \Sigma_j \Big \rVert_\op \leq p^{-2} M \bar \ell \ \text{ with } \ \bar \ell = \frac{1}{n} \sum_{j=1}^{n'} \ell(X'_j). \]
Similarly, we have
\[ \Big \lVert \sum_{j=1}^{n'} \Sigma_j^2 \Big \rVert_\op = \Big( \sum_{j=1}^{n'} \alpha_j^2 \Big) \left \lVert P_{W_\perp} \right \rVert_\op = \frac{1}{n^2 p^2} \sum_{j=1}^{n'} \ell(X'_j)^2 \leq p^{-2} M \bar \ell. \]
Finally,
\[ \max_{1 \leq j \leq n'}\Tr (\Sigma_j) \leq d \max_{1 \leq j \leq n'} \alpha_j = \frac{d}{np} \max_{1 \leq j \leq n'} \ell(X'_j) = p^{-1} M. \]
Gathering the previous inequalities, we finally get
\[ \tilde \E \lVert S - \tilde \E S \rVert_\op \leq \beta C \times \left( p^{-1} (M \bar \ell)^{1/2} + \left( p^{-1} (M \bar \ell)^{1/2} + p^{-1/2} M^{1/2} \right)\log^3(d+n')\right). \]
Recall that $\sup_d \beta < \infty$ since we assume $\inf_d p > 0$. Moreover, Assumption~\ref{ass:MW} implies that $M^{1/2} \log^3(d + n') \Rightarrow 0$. Since $\E(\bar \ell) = 1$, $\bar \ell$ is tight and so the additional multiplicative terms $\bar \ell$ do not change the asymptotic behavior. Finally, we see that the upper bound of the previous display vanishes, which implies that $\tilde \E \lVert S - \tilde \E(S) \rVert_\op \Rightarrow 0$ and so that $\lVert S - \tilde \E(S) \rVert_\op \Rightarrow 0$ as well.


\subsection{Step $3$: $\kappa_* = 1/\lambda_1$ when $V \subset U_\perp$} \label{sub:step-3}

Assume now that $V \subset U_\perp$, and let us prove that $\kappa_* = 1/\lambda_1$. Since $\kappa_* = 1/(1-\gamma_*)$, this amounts to showing that $\gamma_* = 1-\lambda_1$, i.e., in view of~\eqref{eq:MW}, that
\[ \frac{1}{n^\gamma} \max_{1 \leq i \leq n} \xi_A(X_i) \ell(X_i) \Rightarrow \left\{ \begin{array}{ll} 0 & \text{if } \gamma > 1-\lambda_1,\\
	+\infty & \text{if } \gamma < 1-\lambda_1. \end{array} \right. \]
We have
\[ \frac{1}{n^\gamma} \max_{1 \leq i \leq n} \xi_A(X_i) \ell(X_i) = \left( \frac{n'}{qn} \right)^\gamma \times q^\gamma \times \frac{1}{n^{\prime \gamma}} \max_{1 \leq j \leq n'} \ell(P_W X'_j). \]
Since $n'/(qn) \to 1$ in L\textsubscript2 and $\inf_d q > 0$ by Lemma~\ref{lemma:q-n'}, the two first terms of the right-hand side of the previous display do not affect the asymptotic behavior, and we can replace $n'$ by $n$ in the last term. Further,~\eqref{eq:interm} implies that $P_W X'_j$ is equal in distribution to $P_W X$ with $X \sim g$, independently of $n'$. Thus, the result reduces to showing that
\[ \frac{1}{n^\gamma} \max_{1 \leq i \leq n} \ell(X_i) \Rightarrow \left\{ \begin{array}{ll} 0 & \text{if } \gamma > 1-\lambda_1,\\
	+\infty & \text{if } \gamma < 1-\lambda_1. \end{array} \right. \]

We first assume that $\gamma > 1-\lambda_1$, and we prove that $n^{-\gamma} \max \ell(X_i) \Rightarrow 0$. For any $\varepsilon > 0$, we have
\[ \E \left( \left( \frac{1}{n^\gamma} \max_{1 \leq i \leq n} \ell(X_i) \right)^\varepsilon \right) = \frac{1}{n^{\varepsilon \gamma}} \E \left( \max_{1 \leq i \leq n} \ell(X_i)^\varepsilon \right) \leq \frac{1}{n^{\varepsilon \gamma-1}} \E_g \left( \ell(X)^\varepsilon \right). \]
Since $\gamma > 1-\lambda_1$, we can take $1/\gamma < \varepsilon < 1/(1-\lambda_1)$: for this $\varepsilon$, we have $n^{\varepsilon \gamma - 1} \to \infty$ while $\sup_d \E_g(\ell(X)^\varepsilon) < \infty$ by Lemma~\ref{lemma:moments}, which shows that $n^{-\gamma} \max \ell(X_i) \Rightarrow 0$ as desired.

Assume now that $\gamma < 1-\lambda_1$, and let us prove that $n^{-\gamma} \max \ell(X_i) \Rightarrow \infty$. First, write
\[ \sum_{k=1}^r \left( \frac{1}{\lambda_k} - 1 \right) \langle X_i, v_k \rangle^2 = \sum_{k=1}^r \left( 1 - \lambda_k \right) N_{i,k}^2 \]
with $N_{i,k} = \langle v_k, X_i \rangle / \sqrt{\lambda_k}$. Since for $1 \leq k, l \leq r$ with $k \neq l$ we have
\[ \Var_g(v_k^\top X) = v_k^\top \Sigma v_k = \lambda_k \ \text{ and } \ \Cov_g(v_k^\top X, v_l^\top X) = v_k^\top \Sigma v_l = 0 \]
we see that the $N_{i,k}$'s are i.i.d.\ $N(0,1)$ random variables. Moreover, if $Z_i = \sum_{k=1}^r \Indicator{\lambda_k > 1} (\lambda_k - 1) N_{i,k}^2$, we get
\[ \sum_{k=1}^r \left( \frac{1}{\lambda_k} - 1 \right) \langle X_i, v_k \rangle^2 \geq (1-\lambda_1) N_{i,1}^2 - Z_i \]
and so if $I = \arg \max_{1 \leq i \leq n} N_{i,1}^2$, we see starting from~\eqref{eq:expression-ell} that
\[ \max_{1 \leq i \leq n} \ell(X_i) \geq \ell(X_I) = \lvert \Sigma \rvert^{1/2} e^{(1-\lambda_1) N_{I,1}^2 / 2} Z_I. \]
Assumption~\ref{ass:Sigma} implies that $\inf_d \lvert \Sigma \rvert > 0$. Moreover, since the $N_{i,k}$'s are i.i.d., we get that $N_{I,1}$ and $Z_I$ are independent. Thus, in order to show the result, it is enough to prove that
\[ \frac{1}{n^\gamma} e^{(1-\lambda_1) N^2_{I,1} / 2} = \frac{1}{n^\gamma} \max_{1 \leq i \leq n} e^{(1-\lambda_1) N^2_{i,1} / 2} \Rightarrow \infty. \]
Since the $N_{i,1}$ are i.i.d.\ standard Gaussian random variables, it is well-known that $N^2_{I,1} = (\max_{1 \leq i \leq n} N_{i,1})^2$ behaves in first order like $2 \log n$, and so the behavior of $n^{-\gamma} (1-\lambda_1) N^2_{I,1} / 2$ is governed by
\[ \exp \left( - \gamma \log n + (1-\lambda_1) \log n \right) = n^{1-\lambda_1 - \gamma}. \]
Since $\gamma < 1-\lambda_1$, we get as desired $n^{-\gamma} \max \xi_A(X_i) \ell(X_i) \Rightarrow \infty$ when $\gamma < 1 - \lambda_1$.




\subsection{Step $4$: $1 \leq \kappa_* \leq 1/\lambda_1$ when $V \subset U$} \label{sub:step-4}

Assume now that $V \subset U$, and let us prove that $1 \leq \kappa_* \leq 1/\lambda_1$, i.e., $0 \leq \gamma_* \leq 1 - \lambda_1$. That $\gamma_* \geq 0$ is obvious. To prove that $\gamma_* \leq 1-\lambda_1$, we have to prove that
\[ \frac{1}{n^\gamma} \max_{1 \leq i \leq n} \xi_A(X_i) \ell(X_i) \Rightarrow 0 \]
for $\gamma > 1-\lambda_1$. But this readily comes from the inequality
\[ \frac{1}{n^\gamma} \max_{1 \leq i \leq n} \xi_A(X_i) \ell(X_i) \leq \frac{1}{n^\gamma} \max_{1 \leq i \leq n} \ell(X_i) \]
and the fact that $n^{-\gamma} \max_{1 \leq i \leq n} \ell(X_i) \Rightarrow 0$ when $\gamma > 1-\lambda_1$, which we have established above.







\section{Proof of Proposition~\ref{prop:range}}


We now consider the case where $A = \{x: \lvert u^\top x \rvert \leq K\}$ for some $K$ with $\inf_d K > 0$ and $u \in \R^d$ with $\lVert u \rVert = 1$, and where $V = {\rm span}(u) = U$ and any fixed $\lambda_1 \in (0,1)$. It is clear that Assumptions~\ref{ass:FID} and~\ref{ass:Sigma} are satisfied, and that $\inf_d p > 0$.

Let us now consider that $K = 1 + \sqrt{2 \alpha \lambda_1 \log n}$ for some fixed $\alpha \in [0,1]$ and show that Assumption~\ref{ass:MW} is satisfied with $\gamma_* = \alpha (1-\lambda_1)$. In this case, we have
\begin{align*}
	\max_{1 \leq i \leq n} \xi_A(X_i) \ell(X_i) & = \max_{1 \leq i \leq n} \Indicator{\langle X_i, u \rangle^2 \leq K^2} \exp \left( \frac{1}{2} \left( \frac{1}{\lambda_1} - 1 \right) \langle X_i, u \rangle^2 \right)\\
	& = \exp \left( \frac{1}{2} \left( 1 - \lambda_1 \right) \max_{1 \leq i \leq n} \Indicator{N_i^2 \leq K^{\prime 2}} N_i^2 \right)
\end{align*}
with $N_i = \lambda_1^{-1/2} \langle X_i, u \rangle$ and $K' = \lambda_1^{-1/2} K$. Thus, in order to prove the result, it is enough to prove that
\[ \left( 1 - \lambda_1 \right) \max_{1 \leq i \leq n} \Indicator{N_i^2 \leq K^{\prime 2}} N_i^2 - 2 \gamma \log n \Rightarrow \left\{ \begin{array}{ll}
		-\infty & \text{if } \gamma > \alpha (1-\lambda_1),\\
		+\infty & \text{if } \gamma < \alpha (1-\lambda_1),
	\end{array} \right. \]
i.e., that for any $x \in \R$ we have
\begin{equation} \label{eq:goal-A}
	\P \left( (1-\lambda_1) \max_{1 \leq i \leq n} \Indicator{N_i^2 \leq K^{\prime 2}} N_i^2 - 2 \gamma \log n \leq x \right) \to \left\{ \begin{array}{ll}
		1 & \text{if } \gamma > \alpha (1-\lambda_1),\\
		0 & \text{if } \gamma < \alpha (1-\lambda_1).
	\end{array} \right.
\end{equation}
The result for $\alpha = 0$ is obvious since in this case the maximum is bounded (by $K^{\prime 2} = 1$), so assume for the rest of the proof that $\alpha > 0$. Since the $N_i$'s are i.i.d.\ standard Gaussian random variables in dimension one, for any $x \in \R$ we have
\begin{align*}
	\P \Big( (1-\lambda_1) \max_{1 \leq i \leq n} \Indicator{N_i^2 \leq K^{\prime 2}} N_i^2 & - 2 \gamma \log n \leq x \Big)\\
	& = \left( 1 - \P \left( N^2_1 \Indicator{N^2_1 \leq K^{\prime 2}} \geq y \right) \right)^n\\
	& = \left( 1 - \P \left( y \leq N^2_1 \leq K^{\prime 2} \right) \right)^n
\end{align*}
with $y = (1-\lambda_1)^{-1} (x + 2 \gamma \log n)$. Thus,~\eqref{eq:goal-A} is equivalent to
\begin{equation} \label{eq:goal-A-2}
	n \P \left( y \leq N^2_1 \leq K^{\prime 2} \right) \to \left\{ \begin{array}{ll}
		0 & \text{if } \gamma > \alpha (1-\lambda_1),\\
		+\infty & \text{if } \gamma < \alpha (1-\lambda_1).
	\end{array} \right.
\end{equation}
Consider the case $\gamma > \alpha (1-\lambda_1)$. Since $y \sim 2 \gamma (1-\lambda_1)^{-1} \log n$ and $K^{\prime 2} \sim 2 \alpha \log n$, for $n$ large enough we have $y > K^{\prime 2}$ in which case $\P(y \leq N^2_1 \leq K^{\prime 2}) = 0$. Consider now the case $\gamma < \alpha (1-\lambda_1)$. We have
\[ \P \left( N^2_1 \geq y \right) = 2 \P \left( N_1 \geq \sqrt{y} \right) \sim \frac{2 e^{-y/2}}{\sqrt{2\pi y}} \gg \frac{2 e^{-K^{\prime 2}/2}}{\sqrt{2\pi} K'} = \P \left( N^2_1 \geq K^{\prime 2} \right), \]
implying that
\[ n \P \left( y \leq N^2_1 \leq K^{\prime 2} \right) \sim n \P \left( y \leq N^2_1 \right) \sim \sqrt{\frac{2}{\pi}} \exp \left( \log y - y/2 + \log n \right) \]
and since $y/2 \sim \gamma (1-\lambda_1)^{-1} \log n$ with $\gamma (1 - \lambda_1)^{-1} < \alpha \leq 1$, this shows that
\[ n \P \left( y \leq N^2_1 \leq K^{\prime 2} \right) \to \infty. \]
This completes the proof of~\eqref{eq:goal-A-2} and thus of Proposition~\ref{prop:range}.





\appendix

\section{Proof of Theorem~\ref{thm:B-vH-unconditioned}} \label{appendix}

Theorem~\ref{thm:B-vH-unconditioned} is an extension of Theorem 3.21 in~\cite{Brailovskaya24-0}, which we state below.

\begin{thm}[Theorem 3.21 in~\cite{Brailovskaya24-0}]\label{cor:thm3-21}
	Let $S = \sum_{j=1}^n T_j T_j^\top$ where for each $j$, $T_j \sim N(0, \Sigma_j)$ for some covariance matrix $\Sigma_j$ in dimension $d$. Assume that the $T_j$'s are independent. Then there exists a universal, finite constant $C > 0$ such that for any $\varepsilon \in (0,1]$, we have
    \begin{multline*}
		\E \lVert S - \E S \rVert_\op \leq 
        2(1+\varepsilon) \Big \lVert\sum_{i=1}^n \Tr[\Sigma_i] \Sigma_i \Big \rVert_\op^{1/2}\\
		+ \frac{C}{\varepsilon^3} \left(\Big\lVert \sum_{i=1}^n \Sigma_i^2 \Big\rVert_\op^{1/2} + \max_{i\leq n}\Tr \Sigma_i \right)\log^3(d+n).
	\end{multline*}
\end{thm}

We want to apply this result conditionally on $n'$ and the $P_W X'_j$ (thus, under the probability measure $\tilde \P = \P(\cdot \mid n', P_W X'_j, j = 1, \ldots, n')$), and with
\[ T_j = \sqrt{\alpha_j} P_{W_\perp} X'_j \ \text{ with } \ \alpha_j = \frac{1}{np} \ell(P_W X'_j). \]
According to Lemma~\ref{lemma:W}, $P_W X'_j$ (and thus $\alpha_j$) and $P_{W_\perp} X'_j$ are independent. Furthermore, if $V \subset U$, then $X' \sim N(0, P_{W_\perp})$ is a Gaussian vector according to Lemma~\ref{lemma:W}, and so Theorem~\ref{cor:thm3-21} applies with
\[ \Sigma_j = \Var(T_j \mid \alpha_j) = \alpha_j \Var(P_{W_\perp} X') = \alpha_j P_{W_\perp}. \]
This yields the first part of Theorem~\ref{thm:B-vH-unconditioned} when $V \subset U$. When $V \subset U_\perp$, then $P_W X'$ and $P_{W_\perp} X'$ remain independent but $P_{W_\perp} X'$ is no longer Gaussian: Lemma~\ref{lemma:W} shows that it is conditional Gaussian, so Theorem~\ref{cor:thm3-21} does not apply directly and it needs to be extended to this slightly more general case. More precisely, we have to prove the following result.

\begin{prop} \label{prop:extension}
	Let $S = \sum_{j=1}^n T_j T_j^\top$ where for each $j$, $T_j \sim N(0, \Sigma_j) \mid_{A_j}$ for some covariance matrix $\Sigma_j$ in dimension $d$ and some set $A_j$ with $p_j = \P(Z_j \in A_j)$ with $Z_j \sim N(0, \Sigma_j)$. Assume that the $T_j$'s are independent. Then the bound of Theorem~\ref{cor:thm3-21} continues to hold, but with an additive multiplicative factor $1/\min p_j$ in front of the upper bound.
\end{prop}

Before proceeding to the proof of Proposition~\ref{prop:extension}, let us prove that Theorem~$5.9$ is obtained from it.

\begin{proof} [Proof of Theorem~\ref{thm:B-vH-unconditioned} based on Proposition~\ref{prop:extension}]
	The case $V \subset U$ has been discussed above, so assume that $V \subset U_\perp$. We have $T_j = \sqrt{\alpha_j} P_{W_\perp} X'_j$ with $P_{W_\perp} X'_j \sim N(0, P_{W_\perp}) \mid_A$: thus under $\tilde \P$ (in particular, conditionally on $\alpha_j$, independent from $P_{W_\perp} X'_j$), we have $T_j \sim N(0, \alpha_j^{1/4} P_{W_\perp}) \mid_{A_j}$ with $A_j = \{x: \alpha_j^{-1/2} x \in A\}$. In particular, if $Z_j \sim N(0, \alpha_j^{1/4} P_{W_\perp})$ then
	\begin{multline*}
		p_j = \tilde \P \left( Z_j \in A_j \right) = \P \left( \sqrt{\alpha_j} P_{W_\perp} Y \in A_j \mid \alpha_j \right)\\
		= \P \left( P_{W_\perp} Y \in A \mid \alpha_j \right) = \P(P_{W_\perp} Y \in A)
	\end{multline*}
	with $Y \sim N(0,I)$ independent from $\alpha_j$, and using that $P_{W_\perp}^2 = P_{W_\perp}$. Since we are in the case $V \subset U_\perp$, by definition of $W$ we have $W_\perp = U$ and so $p_j = p$ using Assumption~$2$ for the last equality
\end{proof}

In the rest of this document we thus prove Proposition~\ref{prop:extension}, which consists on a slight modification of the proof of~\cite[Theorem 3.21]{Brailovskaya24-0}. Let $\tr = \frac{1}{d} \Tr$. The proof of~\cite[Theorem 3.21]{Brailovskaya24-0} rests on the inequality
\begin{equation} \label{eq:key}
	\E\left[\frac{1}{d}\Tr\left((S - \E S)^{2q}\right) \right]^{1/2q} \leq 2\sigma(S) + C v(S)^{1/2} \sigma(S)^{1/2} q^{3/4} + C R_{2q}(S) q^2
\end{equation}
(in the rest of this document, we adopt the notation of~\cite{Brailovskaya24-0} to which the reader is referred for the various definitions), see~\cite[p.\ 74]{Brailovskaya24-0}. This bound is stated in the Gaussian case but it holds for any distribution of the $T_j$'s. Indeed, if $G$ and $X_\text{free}$ are the Gaussian model and noncommutative model with the same mean and covariance structure as $S$ (see~\cite[sections~$2.1.2$ and~$2.1.3$]{Brailovskaya24-0}), then we have
\begin{align*}
	\E\left[\tr\left((S - \E S)^{2q}\right) \right]^{1/2q} \leq & \left \lvert \E\left[\tr\left((S - \E S)^{2q}\right) \right]^{1/2q} - \E\left[\tr\left(G^{2q}\right) \right]^{1/2q} \right \rvert\\
	& \ \ \ \ \ \ + \left \lvert \E\left[\tr\left(G^{2q}\right) \right]^{1/2q} -  (\tr \otimes 
 \tau)(\lvert X_{\text{free}}\rvert^{2q})^{1/2q} \right \rvert\\
	& \ \ \ \ \ \ + \left \lvert (\tr \otimes 
 \tau)(\lvert X_{\text{free}}\rvert^{2q})^{1/2q} \right \rvert.
\end{align*}
By definition, $G$ and $X_\text{free}$ do not depend on the Gaussian assumption, and so the bounds $2 \sigma(S)$ and $C v(S)^{1/2} \sigma(S)^{1/2} q^{3/4}$ derived in~\cite[Lemma~$2.5$ and Theorem~$2.7$]{Bandeira2023} on the last two terms of the right-hand side of the previous display continue to hold. Only the bound $C R_{2q}(S) q^2$ on the first term $\E\left[\tr\left((S - \E S)^{2q}\right) \right]^{1/2q} - \E\left[\tr\left(G^{2q}\right) \right]^{1/2q}$ depends on the precise nature of $S$, but since this bound is obtained by~\cite[Theorem~$2.9$]{Brailovskaya24-0} which does not depend on the Gaussian assumption, and so we obtain that~\eqref{eq:key} does indeed hold in all generality.

The authors then control the terms $\sigma(S)$, $v(S)$ and $R_q(S)$ in Lemmas~$9.9$ and~$9.10$ of~\cite{Brailovskaya24-0}, where they prove that
\[ \sigma(S) = \Big \lVert\sum_{i=1}^n \left( \Tr[\Sigma_i] \Sigma_i + \Sigma_i^2 \right) \Big \rVert_\op^{1/2}, \ v(S) \leq \sqrt{2} \Big \lVert\sum_{i=1}^n \Sigma_i^2 \Big \rVert_\op^{1/2} \]
and
\[ R_q(S) \leq n^{1/q} \max_{i\leq n} \left( \Tr \Sigma_i + q \lVert \Sigma_i \rVert_\op \right). \]
Thus in order to prove Proposition~\ref{prop:extension}, it is enough to prove that these inequalities continue to hold when the $X'_j$'s are conditional Gaussian, up to an additive multiplicative factor $1/\min p_j$ (and, for $\sigma(S)$, the equality becomes an inequality). In other words, in order to prove Proposition~\ref{prop:extension}, it is enough to prove that, for $S$ as in Proposition~\ref{prop:extension}, we have
\[ \sigma(S) \leq \frac{1}{\min p_j} \Big \lVert\sum_{i=1}^n \left( \Tr[\Sigma_i] \Sigma_i + \Sigma_i^2 \right) \Big \rVert_\op^{1/2}, \ v(S) \leq \frac{\sqrt{2}}{\min p_j} \Big \lVert\sum_{i=1}^n \Sigma_i^2 \Big \rVert_\op^{1/2} \]
and
\[ R_q(S) \leq \frac{1}{\min p_j} n^{1/q} \max_{i\leq n} \left( \Tr \Sigma_i + q \lVert \Sigma_i \rVert_\op \right). \]

We start with the extension for $R_q(S)$ which is straightforward. Following the proof of~\cite[Lemma 9.10]{Brailovskaya24-0}, we have
\[ R_q(S) \leq 2 n^{1/q} \max_{j \leq n} \E \left( \lVert Z_j \rVert_\op^{2q} \mid Z_j \in A_j \right)^{1/q} \leq 2 n^{1/q} \max_{j \leq n} \frac{1}{p_j^{1/q}} \E \left( \lVert Z_j \rVert_\op^{2q} \right)^{1/q} \]
and since $p_j \leq p_j^{1/q}$ since $q \geq 1$, we obtain the result.

Let us now prove the desired bound on $v(S)$. By definition,
\[ v(S)^2 = \sup_{(\Tr \lvert M \rvert)^2 \leq 1} \E \left[  \left( \Tr[M (S - \E S) ]\right)^2 \right]. \]
We have
\begin{align*}
	\E \left\{ \left[ \Tr \left( M (S - \E S) \right) \right]^2 \right\} &= \E \left\{ \left[ \Tr \left( \sum_{j=1}^n (M T_j T_j^\top - \E(M T_j T_j^\top) \right) \right]^2 \right\}\\
	&= \E \left\{ \left[ \sum_{j=1}^n \left( \Tr  (M T_j T_j^\top - \E(M T_j T_j^\top) \right) \right]^2 \right\}\\
	&= \sum_{j=1}^n \E \left[ \left( \Tr  (M T_j T_j^\top) - \E(\Tr(M T_j T_j^\top)) \right)^2 \right]
\end{align*}
using the definition of $S$ for the first equality, the linearity of the trace for the second one, and the fact that the $T_j$ are independent, so that the variance of the sum is the sum of variance, for the third one. If $Z_j \sim N(0, \Sigma_j)$, we have
\begin{multline*}
	\E \left[ \left( \Tr  (M T_j T_j^\top) - \E(\Tr(M T_j T_j^\top)) \right)^2 \right]\\
	= \E \left[ \left( \Tr  (M Z_j Z_j^\top) - \E(\Tr(M Z_j Z_j^\top) \mid Z_j \in A_j) \right)^2 \mid Z_j \in A_j \right]
\end{multline*}
and so the variational characterization of the mean gives that
\begin{multline*}
	\E \left[ \left( \Tr  (M T_j T_j^\top) - \E(\Tr(M T_j T_j^\top)) \right)^2 \right]\\
	\leq \E \left[ \left( \Tr  (M Z_j Z_j^\top) - \E(\Tr(M Z_j Z_j^\top)) \right)^2 \mid Z_j \in A_j \right]
\end{multline*}
which leads to
\[ \E \left[ \left( \Tr  (M T_j T_j^\top) - \E(\Tr(M T_j T_j^\top)) \right)^2 \right] \leq \frac{1}{p_j} \E \left[ \left( \Tr  (M Z_j Z_j^\top) - \E(\Tr(M Z_j Z_j^\top)) \right)^2 \right]. \]
This bound allows to resume the original proof in the Gaussian setting, with no conditioning, at the expense of an additional multiplicative $1/\min p_j$ term.

Finally, the bound on $\sigma(S)$ proceeds with the same arguments as for $v(S)$, but in a multi-dimensional case. Indeed, we have
\begin{align*}
	\sigma(S)^2 & = \left \lVert \E \left[ \left( S - \E(S) \right)^2 \right] \right \rVert_\op\\
	& = \left \lVert \sum_{i=1}^n \E \left[ \left( T_j T_j^\top - \E(T_j T_j^\top) \right)^2 \right] \right \rVert_\op\\
	& = \left \lVert \sum_{i=1}^n \E \left[ \left( Z_j Z_j^\top - \E(Z_j Z_j^\top \mid Z_j \in A_j) \right)^2 \mid Z_j \in A_j \right] \right \rVert_\op.
\end{align*}
The following lemma is a multi-dimensional generalization of the inequality used previously to control $v(S)$.

\begin{lemma}
	For any random matrix $M$ and any event $E$, we have
	\[ \E \left[ \left( M - \E(M \mid E)\right)^2 \mid E \right] \leq \frac{1}{\P(E)} \E \left[ \left( M - \E(M)\right)^2 \right] \]
	where $A \leq B$ means that $B-A$ is positive.
\end{lemma}

\begin{proof}
	We have
	\[ \E \left[ \left( M - \E(M)\right)^2 \mid E \right] - \E \left[ \left( M - \E(M \mid E)\right)^2 \mid E \right] = \left[ \E(M) - \E(M \mid E) \right]^2 \]
	implying that $\E \left[ \left( M - \E(M \mid E)\right)^2 \mid E \right] \leq \E \left[ \left( M - \E(M)\right)^2 \mid E \right]$.
\end{proof}

Thus
\begin{multline*}
	\E \left[ \left( Z_j Z_j^\top - \E(Z_j Z_j^\top \mid Z_j \in A_j) \right)^2 \mid Z_j \in A_j \right]\\
	\leq \frac{1}{\min p_j} \E \left[ \left( Z_j Z_j^\top - \E(Z_j Z_j^\top) \right)^2 \right]
\end{multline*}
and so
\begin{multline*}
	\sum_j \E \left[ \left( Z_j Z_j^\top - \E(Z_j Z_j^\top \mid Z_j \in A_j) \right)^2 \mid Z_j \in A_j \right]\\
	\leq \frac{1}{\min p_j} \sum_j \E \left[ \left( Z_j Z_j^\top - \E(Z_j Z_j^\top) \right)^2 \right].
\end{multline*}
Since $\lVert A \rVert_\op \leq \lVert B \rVert_\op$ when $A,B$ are positive with $A \leq B$, we obtain
\[ \sigma(S)^2 \leq \frac{1}{\min p_j} \left \lVert \sum_{i=1}^n \E \left[ \left( Z_j Z_j^\top - \E(Z_j Z_j^\top) \right)^2 \right] \right \rVert_\op \]
and so we are back to the Gaussian setting.


\begin{thebibliography}{10}

\bibitem{agapiou_importance_2017}
S.~Agapiou, O.~Papaspiliopoulos, D.~Sanz-Alonso, and A.~M. Stuart.
\newblock Importance sampling: Intrinsic dimension and computational cost.
\newblock {\em Statistical Science}, 32(3):405--431, 2017.

\bibitem{BAIK20061382}
Jinho Baik and Jack~W. Silverstein.
\newblock Eigenvalues of large sample covariance matrices of spiked population
  models.
\newblock {\em Journal of Multivariate Analysis}, 97(6):1382--1408, 2006.

\bibitem{Bandeira2023}
Afonso~S. Bandeira, March~T. Boedihardjo, and Ramon van Handel.
\newblock Matrix concentration inequalities and free probability.
\newblock {\em Inventiones mathematicae}, 234(1):419--487, Oct 2023.

\bibitem{bassamboo_portfolio_2008}
Achal Bassamboo, Sandeep Juneja, and Assaf Zeevi.
\newblock Portfolio {Credit} {Risk} with {Extremal} {Dependence}: {Asymptotic}
  {Analysis} and {Efficient} {Simulation}.
\newblock {\em Operations Research}, 56(3):593--606, June 2008.

\bibitem{beh2023insight}
Jason Beh, Yonatan Shadmi, and Florian Simatos.
\newblock {Insight from the Kullback--Leibler divergence into adaptive
  importance sampling schemes for rare event analysis in high dimension}.
\newblock {\em The Annals of Applied Probability}, 35(2):1083--1124, 2025.

\bibitem{benaych-georges_eigenvalues_2011}
Florent Benaych-Georges and Raj~Rao Nadakuditi.
\newblock The eigenvalues and eigenvectors of finite, low rank perturbations of
  large random matrices.
\newblock {\em Advances in Mathematics}, 227(1):494--521, May 2011.

\bibitem{Brailovskaya24-0}
Tatiana Brailovskaya and Ramon van Handel.
\newblock Universality and sharp matrix concentration inequalities.
\newblock {\em Geometric and Functional Analysis}, 34(6):1734--1838, Dec 2024.

\bibitem{capitaine2016spectrumdeformedrandommatrices}
Mireille Capitaine and Catherine Donati-Martin.
\newblock Spectrum of deformed random matrices and free probability.
\newblock In {\em Advanced topics in random matrices}, volume~53 of {\em Panor.
  Synth\`eses}, pages 151--190. Soc. Math. France, Paris, 2017.

\bibitem{chan_improved_2012}
Joshua C.~C. Chan and Dirk~P. Kroese.
\newblock Improved cross-entropy method for estimation.
\newblock {\em Statistics and Computing}, 22(5):1031--1040, September 2012.

\bibitem{Paul_07}
Paul Debashis.
\newblock Asymptotics of sample eigenstructure for a large dimensional spiked
  covariance model.
\newblock {\em Statistica Sinica}, 17(4):1617--1642, 2007.

\bibitem{Durrett_2019}
Rick Durrett.
\newblock {\em Probability: Theory and Examples}.
\newblock Cambridge Series in Statistical and Probabilistic Mathematics.
  Cambridge University Press, 5 edition, 2019.

\bibitem{ELMASRI2021}
Maxime {El Masri}, J{\'e}r{\^o}me Morio, and Florian Simatos.
\newblock Improvement of the cross-entropy method in high dimension for failure
  probability estimation through a one-dimensional projection without gradient
  estimation.
\newblock {\em Reliability Engineering \& System Safety}, 216:107991, 2021.

\bibitem{el_masri2024}
Maxime El~Masri, J{\'e}r{\^o}me Morio, and Florian Simatos.
\newblock Optimal projection for parametric importance sampling in high
  dimensions.
\newblock {\em Computo}, 2024.

\bibitem{lam2019}
Clifford Lam.
\newblock High-dimensional covariance matrix estimation.
\newblock {\em WIREs Computational Statistics}, 12(2):e1485, 2020.

\bibitem{papaioannou_improved_2019}
Iason Papaioannou, Sebastian Geyer, and Daniel Straub.
\newblock Improved cross entropy-based importance sampling with a flexible
  mixture model.
\newblock {\em Reliability Engineering \& System Safety}, 191:106564, November
  2019.

\bibitem{papaioannou_sequential_2016}
Iason Papaioannou, Costas Papadimitriou, and Daniel Straub.
\newblock Sequential importance sampling for structural reliability analysis.
\newblock {\em Structural Safety}, 62:66--75, 2016.

\bibitem{PAPAKONSTANTINOU2023}
Konstantinos~G. Papakonstantinou, Hamed Nikbakht, and Elsayed Eshra.
\newblock Hamiltonian mcmc methods for estimating rare events probabilities in
  high-dimensional problems.
\newblock {\em Probabilistic Engineering Mechanics}, 74:103485, 2023.

\bibitem{uribe2021}
Felipe Uribe, Iason Papaioannou, Youssef~M. Marzouk, and Daniel Straub.
\newblock Cross-entropy-based importance sampling with failure-informed
  dimension reduction for rare event simulation.
\newblock {\em SIAM/ASA Journal on Uncertainty Quantification}, 9(2):818--847,
  2021.

\end{thebibliography}

\end{document}


